\documentclass[11pt,a4paper]{article}
\usepackage[T1]{fontenc}
\usepackage[utf8]{inputenc}
\usepackage{lmodern}
\usepackage[margin=2.2cm]{geometry}
\usepackage{amsmath,amssymb}
\usepackage{listings}
\usepackage{xcolor}
\usepackage{booktabs}
\usepackage{enumitem}
\usepackage{hyperref}
\hypersetup{colorlinks,linkcolor=black,urlcolor=blue!50!black}

% notation of the TEQ paper
\newcommand{\bl}[1]{\lfloor #1\rfloor}                       % butlast
\newcommand{\tr}[3]{\mathsf{trace}(#1,\ \mathsf{for}\ #2\ \mathsf{do}\ #3)}
\newcommand{\mapc}[2]{\mathsf{map}(\mathsf{for}\ #1\ \mathsf{yield}\ #2)}
\newcommand{\al}[2]{\mathsf{all}(\mathsf{for}\ #1\ \mathsf{check}\ #2)}
\newcommand{\cnt}{\mathsf{counter}}
\newcommand{\cur}{\mathsf{current}}
\newcommand{\intf}{\mathsf{integer}}
\newcommand{\natf}{\mathsf{natural}}
\newcommand{\sums}{\mathsf{sums}}
\newcommand{\sorted}{\mathsf{sorted}}
\newcommand{\perm}{\mathsf{perm}}
\newcommand{\partitioned}{\mathsf{partitioned}}
\newcommand{\argmin}{\mathsf{argmin}}
\newcommand{\swap}[3]{\mathsf{swap}(#1, #2, #3)}
\newcommand{\ite}[3]{\mathsf{if}\ #1\ \mathsf{then}\ #2\ \mathsf{else}\ #3}
\newcommand{\upd}[2]{\mathsf{case}\ #1 \Rightarrow #2}
\newcommand{\lemma}[1]{\textsf{#1}}
\newcommand{\seq}[1]{\hat{#1}}                               % the fresh sequences of inner loops

\lstdefinelanguage{dafny}{
  morekeywords={method,returns,requires,ensures,var,while,if,else,function,lemma,forall,int,bool,array},
  sensitive=true, morecomment=[l]{//}, morestring=[b]",
}
\lstset{language=dafny,basicstyle=\ttfamily\footnotesize,keywordstyle=\bfseries,commentstyle=\itshape\color{black!60},
  columns=fullflexible,keepspaces=true,xleftmargin=1.5em}

\title{Trace Equation Systems for Array Algorithms:\\ Worked Examples}
\author{Gidon Ernst\\[2pt] \small LMU Munich, Germany\\ \small\texttt{gidon.ernst@gmail.com}}
\date{October 2026}

\everydisplay{\small}
\begin{document}
\maketitle

\noindent
These notes document the examples of the folder \texttt{examples/dafny/algorithms} of Cuvée in the style of the
paper \emph{Trace Equation Systems for Program Verification}: for each program we show the trace
equation system that Cuvée generates (\texttt{-teq}), the lemmas that the simplifier applies
(\texttt{-simplify}), how they are instantiated, and the conceptual idea behind them. All lemmas are
stated in the Dafny dialect of Cuvée at the beginning of the respective file; the lemma database (the
modules of \texttt{resources/cuvee/teq}) is injected by \texttt{-teq}, and the module of the fragments of
inner loops by \texttt{-include teq:nested} where a file uses it. We write $\bl{X}$ for \texttt{butlast(X)}, the trace $X$
without its last element, $\cur(X)$ for its last element, and $\cnt(i, c, n)$ for the counter
$\langle i, i + c, \ldots, i + n c\rangle$ of $n$ steps. The sections are ordered by the complexity of the proof, from a
single loop with two lemmas to loops in sequence and nested loops with summaries and inductions;
Section~\ref{sec:lb} documents an attempt that is open. Sections~\ref{sec:monocera} to~\ref{sec:ifeqn2}
re-trace the proofs of some of the C programs of the MonoCera benchmarks with the database alone, and
Section~\ref{sec:principles} collects the principles that the examples share, and Section~\ref{sec:conclusion}
reads them as symptoms of the engine and draws the conclusion for the representation. Capital letters $K, I, B, M$ denote traces,
hatted letters $\seq{M}, \seq{J}, \seq{N}$ the fresh \emph{sequences} that the summary of an inner loop
introduces, one entry per iteration of the outer loop. $\swap{x}{p}{q}$ abbreviates the double store
$x[p := x[q]][q := x[p]]$, which is how the swaps appear in the trace equations.

\section{Linear Search: an Iteration Constraint Abstracted and Instantiated}\label{sec:lsearch}

\subsection{The program}
\begin{lstlisting}
method lsearch(a: array<int, int>, n: int, x: int) returns (i: int)
  requires 0 <= n
  ensures 0 <= i <= n
  ensures i < n ==> a[i] == x
  ensures forall k: int :: 0 <= k < i ==> a[k] != x
{
  i := 0;
  while i < n && a[i] != x { i := i + 1; }
}
\end{lstlisting}
The first index of $x$ among $a[0..n)$, or $n$. The loop test is data-dependent, so the number of
iterations is not determined by $n$; the trace equation system records the test as the iteration
constraint and its negation as the exit.

\subsection{The trace equation system}
\begin{align}
  K &= \cnt(0, 1, n') \qquad i = \cur(I) \qquad I = \tr{0}{\bl{K}}{i, k \mapsto i + 1} \label{eq:ls0}\\
  &\al{\bl{K}, \bl{I}}{k, i \mapsto i < n \wedge a[i] \ne x} \label{eq:ls1}\\
  &\neg(i < n \wedge a[i] \ne x) \label{eq:ls2}\\
  &\Longrightarrow\ 0 \le i \le n\ \wedge\ (i < n \Rightarrow a[i] = x)\ \wedge\ (\forall k.\ 0 \le k < i \Rightarrow a[k] \ne x) \nonumber
\end{align}
With \texttt{-split}, the conclusion becomes four goals, one per conjunct, the last two with their premises
as hypotheses: $i < n \Rightarrow a[i] = x$, and $0 \le k < i \wedge a[k] = x \Rightarrow \mathsf{false}$.
The simplifier's own rules do the following: $I$ is a counter by \lemma{counter\_def}, so
$i = \intf(n')$; the constraint \eqref{eq:ls1} is distributed into
$\al{\bl{\cnt(0, 1, n')}}{k \mapsto k < n}$, which \lemma{invs\_counter\_upper} resolves to
$n' = 0 \vee \intf(n') \le n$ (hence $0 \le i \le n$, the first two goals), and
$\al{\bl{\cnt(0, 1, n')}}{k \mapsto a[k] \ne x}$, the search constraint proper; the exit
\eqref{eq:ls2} is $n \le i \vee a[i] = x$, which with $i < n$ gives the third goal.

\subsection{The lemmas}
The last goal needs the search constraint at a witness: for $0 \le k < i$, the constraint says
$a[k] \ne x$. The paper's lemma database has no instantiation of iteration constraints at positions
(\texttt{TODO.md}: no instantiation at a witness); the two lemmas of \texttt{lsearch.dfy}, since merged
into \texttt{teq.dfy}, abstract the constraint by an atom and instantiate the atom:
\begin{align*}
  &\lemma{not\_occurs\_selects\_def}:\\
  &\quad \al{\bl{\cnt(i, 1, n)}}{k \mapsto a[k] \ne x}\ \Longleftrightarrow\ \neg\mathsf{occurs}(x, \mathsf{selects}(a, \bl{\cnt(i, 1, n)}))\\[4pt]
  &\lemma{select\_not\_occurs\_selects}:\\
  &\quad \neg\mathsf{occurs}(x, \mathsf{selects}(a, \bl{\cnt(i, 1, n)}))\ \wedge\ i \le k < i + \intf(n)
    \ \Longrightarrow\ (a[k] = x) \Longleftrightarrow \mathsf{false}
\end{align*}
The first matches the constraint with $i := 0$, $n := n'$ (the counter $K$ is followed in the dependency
position) and replaces it by the atom "$x$ does not occur among the values of $a$ at the positions
$0, \ldots, n' - 1$". The second is a rule on the atom $a[k] = x$ of the last goal: its side conditions
are the fact just produced, matched with $i := 0$, $n := n'$, and the bounds $0 \le k < \intf(n')$,
decided from $0 \le k < i$ and $i = \intf(n')$. The hypothesis $a[k] = x$ becomes $\mathsf{false}$ and
the goal closes.

This is the pattern of \lemma{not\_occurs\_prefix\_def} and \lemma{earlier\_not\_occurs\_prefix} of
\texttt{teq.dfy} for reads $\mathsf{earlier}(k, xs)$ of a sequence, here for reads of an array: an
iteration constraint becomes a named fact about the prefix visited, and a second lemma uses the fact at
any position in the prefix. \texttt{TODO.md} notes that a side condition that is the iteration constraint
itself would make the rule cycle (finding 19 of \texttt{NOTES.md}), which is why the abstraction by an
atom is not merely convenient.

The file's header prescribes \texttt{-teq -split -simplify}, since without the split the disjunction of
the two exits, $n \le i \vee a[i] = x$, had to be split by the simplifier against the postcondition; with
the current version the goal closes without \texttt{-split} as well.

\section{Nested Sum: the Fragment Lemmas Alone}\label{sec:nsum}

\subsection{The program}
\begin{lstlisting}
method main(m: int, n: int) returns (c: int)
  requires 0 <= m && 0 <= n
  ensures c == m * n
{
  var i: int; var j: int;
  c := 0; i := 0;
  while i < m {
    j := 0;
    while j < n { j := j + 1; }
    c := c + j;
    i := i + 1;
  }
}
\end{lstlisting}
The smallest nested loop whose result depends on the inner loop: $c$ sums the final values of $j$,
which are all $n$. There is nothing to summarize but the inner counter, so this example introduces the
\emph{fragment lemmas} of \texttt{fragments.dfy}, which every nested example uses (Section~\ref{sec:sel:frag}
states the first two in full), and shows how an iteration constraint reaches a lemma of \texttt{teq.dfy}.

\subsection{The trace equation system}
\begin{align}
  K &= \cnt(0, 1, n') \qquad I = \cnt(0, 1, n') \qquad \intf(n') = m \qquad c = \mathsf{sum}(\seq{J}) \label{eq:ns0}\\
  &\al{\bl{K}, \bl{I}, \seq{N}, \seq{K}'}{k, i, n'', K' \mapsto K' = \cnt(0, 1, n'')} \label{eq:ns1}\\
  &\al{\bl{K}, \bl{I}, \seq{K}', \seq{J}'}{k, i, K', J \mapsto J = \sums(0, \mapc{\bl{K'}}{1})} \label{eq:ns2}\\
  &\al{\bl{K}, \bl{I}, \seq{J}, \seq{J}'}{k, i, j, J \mapsto j = \cur(J)} \label{eq:ns3}\\
  &\al{\bl{K}, \bl{I}, \seq{J}'}{k, i, J \mapsto \al{\bl{J}}{j \mapsto j < n}} \label{eq:ns4}\\
  &\al{\bl{K}, \bl{I}, \seq{J}}{k, i, j \mapsto n \le j} \label{eq:ns5}\\
  &\Longrightarrow\ c = m \cdot n \nonumber
\end{align}
The trace of $c$ is $\tr{0}{\bl{K}, \seq{J}}{c, k, j \mapsto c + j}$, which the introduction rule of
$\sums$ in \texttt{teq.dfy} folds and the acceleration rule $\cur(\sums(x, V)) = x + \mathsf{sum}(V)$
turns into $c = \mathsf{sum}(\seq{J})$: the sum of the final values of the inner loop. The task is to
show that $\seq{J}$ consists of $n$ repeated $m$ times.

\subsection{The lemmas and their instantiation}
Two lemmas of \texttt{fragments.dfy}, stated in full in Section~\ref{sec:sel:frag}, apply with
$X := \lambda k.\, 0$ and $c := 1$: \lemma{invs\_sums\_counter} (the trace of an inner variable increased
by $c$ over the inner counter $K'$, which is $\cnt(0, 1, n'')$ in the same iteration, is a counter from
its initial value $X(k)$) turns \eqref{eq:ns2} into
$\al{\bl{K}, \seq{N}, \seq{J}'}{k, n'', J \mapsto J = \cnt(0, 1, n'')}$, and
\lemma{invs\_current\_counter} (the final value of an inner counter) turns \eqref{eq:ns3} into
$\al{\bl{K}, \seq{N}, \seq{J}}{k, n'', j \mapsto j = \intf(n'')}$. Two more fragment lemmas of
\texttt{fragments.dfy} resolve the test and the exit; both are stated for an inner counter from $0$,
since their side conditions must be in the normal form of the facts, where a general initial value
$X(k) + \intf(n'')$ would not match $\intf(n'')$ (a higher-order summand matches a single summand):
\begin{align*}
  &\lemma{invs\_invs\_counter\_upper\_zero}:\\
  &\quad 0 \le n\ \wedge\ \al{ks, Ns, Js}{k, n'', J \mapsto J = \cnt(0, 1, n'')}\\
  &\quad\Longrightarrow\\
  &\quad \al{ks, Js}{k, J \mapsto \al{\bl{J}}{x \mapsto x < n}}\\
  &\quad\Longleftrightarrow\ \al{ks, Ns}{k, n'' \mapsto \intf(n'') \le n}\\[4pt]
  &\lemma{invs\_exit\_counter\_zero}:\\
  &\quad \al{ks, Ns, js}{k, n'', j \mapsto j = \intf(n'')}\\
  &\quad\wedge\ \al{ks, Ns}{k, n'' \mapsto \intf(n'') \le n}\ \Longrightarrow\\
  &\quad \al{ks, js}{k, j \mapsto n \le j}\ \Longleftrightarrow\ \al{ks, js}{k, j \mapsto j = n}
\end{align*}
The first matches the test \eqref{eq:ns4}, with the counter fragment as its side condition (which binds
$Ns := \seq{N}$); the inner constraint $j < n$ over $\bl{\cnt(0, 1, n'')}$ is the lemma
\lemma{invs\_counter\_upper} of \texttt{teq.dfy}, $n'' = 0 \vee (0 < n \wedge \intf(n'') - 1 < n)$,
which the nonnegative bound simplifies to $\intf(n'') \le n$. The second matches the exit
\eqref{eq:ns5}, finds the two facts, and states the final value: $j = n$ in each outer iteration. It is
an equivalence under its side conditions, since $j = \intf(n'') \le n$.

The goal is now $\mathsf{sum}(\seq{J}) = m \cdot n$ with the fact
$\al{\bl{K}, \seq{J}}{k, j \mapsto j = n}$. In the body of that iteration constraint the simplifier
reads $j$ as $\mathsf{earlier}(\natf(k), \seq{J})$, so the fact is
\[ \al{\bl{\cnt(0, 1, n')}}{k \mapsto \mathsf{earlier}(\natf(k), \seq{J}) = n} \]
(the counter $K$ is followed in the dependency position), which \lemma{consts\_prefix\_def} of \texttt{teq.dfy} abstracts
by the atom $\mathsf{consts\_prefix}(n', n, \seq{J})$, and \lemma{sum\_consts\_prefix} with that atom as
its side condition rewrites $\mathsf{sum}(\seq{J})$ to $n \cdot \intf(n')$. (Until October 10 this
was \lemma{sum\_invs\_const}, with the iteration constraint itself as the side condition, which cycles
in general, see \texttt{TODO.md}.) The goal $\intf(n') \cdot n = m \cdot n$ with $\intf(n') = m$ is
closed by a lemma \lemma{mul\_eq\_cancel}, $(a \cdot n = b \cdot n) \Leftrightarrow (n = 0 \vee a = b)$,
since the simplifier substitutes the solution of $\intf(n') = m$ into polynomials but not into the
product.

\subsection{The idea}
The fragments of an inner loop are a trace equation system over bound variables. The simplifier would
resolve them if the bound variables of one fragment saw the equations of the others, which they do not
(\texttt{NOTES.md}, the two problems of nested loops). The fragment lemmas are the four steps of that
resolution as rewrite rules: the inner trace is a counter, its current value is its length, the test
bounds the length, the exit fixes it. The same four lemmas serve every example with an inner counter
from $0$ (\texttt{nested\_mult}, \texttt{row\_sums}, \texttt{transpose}); with another initial value
the test and the exit are instead kept as side conditions of a summary (Section~\ref{sec:sel:sum}).

\section{Row Sums and Transpose: a Chain of Specific Lemmas}\label{sec:rows}

\subsection{The program}
\begin{lstlisting}
method main(A: array<int, array<int, int>>, m: int, n: int) returns (s: array<int, int>)
  requires 0 <= m && 0 <= n
  ensures forall i: int :: 0 <= i < m ==>
    s[i] == sum(selects(A[i], butlast(counter(0, 1, natural(n)))))
{
  var i: int;
  var j: int;
  var t: int;
  i := 0;
  while i < m
  {
    t := 0;
    j := 0;
    while j < n
    {
      t := t + A[i][j];
      j := j + 1;
    }
    s := s[i := t];
    i := i + 1;
  }
}
\end{lstlisting}
The inner loop sums one row and the outer loop stores the sum at the position of its iteration; the
transpose (\texttt{transpose.dfy}) is the same shape with an array-valued inner result, the stores of
column $k$ into $B$. Both are run with \texttt{-include teq:nested}, the module of the fragments of inner loops.

\subsection{The trace equation system}
With the module, the goal of the program alone is
\begin{align}
  K &= I = \cnt(0, 1, n') \qquad \intf(n') = m \label{eq:rs0}\\
  &\al{\bl{K}, \bl{I}, \seq{N}, \seq{K}'}{k, i, nn, K' \mapsto K' = \cnt(0, 1, nn)} \nonumber\\
  &\al{\bl{K}, \seq{N}, \seq{J}'}{k, nn, J \mapsto J = \cnt(0, 1, nn)} \label{eq:rs1}\\
  &\seq{N} = \bl{\mathsf{consts}(n', \natf(n))} \qquad \mathsf{consts\_prefix}(n', n, \seq{J}) \label{eq:rs2}\\
  &\al{\bl{K}, \bl{I}, \seq{T}', \seq{J}'}{k, i, T, J \mapsto T = \sums(0, \mathsf{selects}(A[k], \bl{J}))} \label{eq:rs3}\\
  &\al{\bl{K}, \bl{I}, \seq{T}, \seq{T}'}{k, i, t, T \mapsto t = \cur(T)} \label{eq:rs4}\\
  &\al{\bl{K}, \bl{I}, \seq{J}}{k, i, j \mapsto n \le j} \qquad \ldots \nonumber\\
  &\Longrightarrow\ 0 \le i < m \Rightarrow \mathsf{earlier}(\natf(i), \seq{T}) = \mathsf{sum}(\mathsf{selects}(A[i], \bl{\cnt(0, 1, \natf(n))})) \nonumber
\end{align}
The module has resolved the inner counter \eqref{eq:rs1}, the exit (abstracted to $\mathsf{consts\_prefix}$)
and the numbers of inner iterations \eqref{eq:rs2}: $\seq{N}$ is the constant sequence $n$. The sum of
row $k$ is the trace \eqref{eq:rs3} over the inner counter, its result the fragment \eqref{eq:rs4}, and the
postcondition reads the stored results at position $i$. What is left is to carry $\seq{N}$ into the counter,
the counter into the sum, the sum into the result, and the results into the read: four substitutions
between fragments over different dependency lists, which the simplifier does not make by itself.

\subsection{The chain}
The first version of this file stated the substitutions as general lemmas: a fragment about the inner
counter $J$ with a higher-order body $t = F(k, J)$ is stated about $\cnt(0, 1, nn)$, and a read
$\mathsf{earlier}(\natf(i), ts)$ is the body of the $\mathsf{all}$ over $ts$ at $i$. Both have a general
left-hand side, matching every fragment, respectively every read, together with an $\mathsf{all}$ side
condition over the lists of the left-hand side; matching the condition builds the simplifier's context of
the fact's combinator, where the reads of the dependencies are rewritten, on which the rule matches again,
one level deeper, until the step budget is exhausted (finding 19 of \texttt{NOTES.md}, the cycle that
\texttt{TODO.md} describes for \lemma{sum\_invs\_const}). The runs did not terminate, which the runners of
the time reported as proofs, since a killed run prints no goal.

The version that verifies has one specific lemma per step, each with the fragment it depends on as a side
condition over a fresh list variable, identified by an equation (Section~\ref{sec:cock}); $C_n$ abbreviates
the counter $\bl{\cnt(0, 1, \natf(n))}$:
\begin{align*}
  &\lemma{invs\_counter\_consts}:\quad 0 \le n\\
  &\quad\Longrightarrow\ \al{ks, \bl{\mathsf{consts}(m, \natf(n))}, Js}{k, nn, J \mapsto J = \cnt(0, 1, nn)}\\
  &\qquad\Longleftrightarrow\ \al{ks, Js}{k, J \mapsto \bl{J} = C_n}\\[2pt]
  &\lemma{invs\_sums\_selects\_counter}:\quad \al{ks, Js_2}{k, J \mapsto \bl{J} = C_n}\ \wedge\ Js_2 = Js\\
  &\quad\Longrightarrow\ \al{ks, Ts, Js}{k, T, J \mapsto T = \sums(0, \mathsf{selects}(F(k), \bl{J}))}\\
  &\qquad\Longleftrightarrow\ \al{ks, Ts}{k, T \mapsto T = \sums(0, \mathsf{selects}(F(k), C_n))}\\[2pt]
  &\lemma{invs\_current\_def}:\quad \al{ks, Ts_2}{k, T \mapsto T = G(k)}\ \wedge\ Ts_2 = Ts\\
  &\quad\Longrightarrow\ \al{ks, ts, Ts}{k, t, T \mapsto t = \cur(T)}\\
  &\qquad\Longleftrightarrow\ \al{ks, ts}{k, t \mapsto t = \cur(G(k))}\\[2pt]
  &\lemma{invs\_maps\_sum}:\quad \al{ks, ts}{k, t \mapsto t = \mathsf{sum}(F(k))}\\
  &\qquad\Longleftrightarrow\ ts = \mapc{ks}{k \mapsto \mathsf{sum}(F(k))}\\[2pt]
  &\lemma{earlier\_maps\_counter}:\quad 0 \le i < \intf(m)\\
  &\quad\Longrightarrow\ \mathsf{earlier}(\natf(i), \mapc{\bl{\cnt(0, 1, m)}}{k \mapsto F(k)}) = F(i)
\end{align*}
The first matches \eqref{eq:rs1} with the dependency $\bl{\mathsf{consts}(m, \natf(n))}$ as a pattern,
which is matched against $\seq{N}$ through its equation \eqref{eq:rs2} in the context (a variant with
$\seq{N}$ as a variable and the equation as a side condition is in the file as well). The second rewrites
\eqref{eq:rs3} to a sum over the counter $\bl{\cnt(0, 1, \natf(n))}$, with $F := \lambda k.\ A[k]$. The third
rewrites \eqref{eq:rs4} with $G$ the sum of row $k$, after which the database makes $\cur(\sums(0, xs))$ into
$\mathsf{sum}(xs)$; the fourth turns the pointwise fragment of the results
into a map (specific to sums, so that the fragments $j = \cur(J)$ of counters are left to the module); the
fifth reads the map at $i$, with $m := n'$ and $0 \le i < \intf(n')$ decided from $0 \le i < m$ and
\eqref{eq:rs0}. Every left-hand side is a fragment with a function symbol on the side of its bound variable
or a read of a $\mathsf{map}$, and every $\mathsf{all}$ condition is over a fresh list: nothing matches what
the rewriting of a read in a combinator context produces, and the runs take under twenty seconds.

For the transpose the chain is the same with the inner stores: the trace of stores over $\bl{J}$ is folded
into $\mathsf{column}(h, A, k, \natf(n))$ once $J = \cnt(0, 1, \natf(n))$ is known, the final values of the
columns are the fragment $b = \cur(\mathsf{column}(h, A, k, \natf(n)))$ over the previous value $h$, then
\lemma{current\_results\_steps} (Section~\ref{sec:nsum}) makes the outer trace the trace of the columns, and
\lemma{current\_transpose} reads the transposed element.

\subsection{The idea}
An inner result stored at the position of the outer iteration is the pattern of the paper's nested loops,
and with the module it needs no new concept, only discipline: every step is a lemma whose left-hand side is
specific and whose side condition is a fragment over a fresh list, matched against the facts. The general
substitution lemmas that one would write first are exactly the rules that the cycle mechanism of
\texttt{TODO.md} forbids; the chain makes the same substitutions one fragment at a time.

\section{Selection Sort: Fragments, a Summary, and an Induction}\label{sec:sel}

\subsection{The program}

\begin{lstlisting}
method main(a: array<int, int>, n: int) returns (b: array<int, int>)
  requires 0 <= n
  ensures sorted(b, 0, n)
  ensures perm(b, a)
{
  var i: int; var j: int; var m: int; var t: int;
  b := a; i := 0;
  while i < n {
    m := i; j := i + 1;
    while j < n {
      if b[j] < b[m] { m := j; }
      j := j + 1;
    }
    t := b[i]; b[i] := b[m]; b[m] := t;
    i := i + 1;
  }
}
\end{lstlisting}

\noindent
Arrays are total maps $\mathsf{array}\langle\mathsf{int},\mathsf{int}\rangle$; $\sorted(a, lo, hi)$
states that $a[lo..hi)$ is nondecreasing, and $\perm(a, b)$ that $a$ is a permutation of $b$ as a total
map, i.e. $b$ composed with a bijection of the indices. Both are abstract predicates, declared without
a definition; their properties are the lemmas below. Whole-array permutation needs no side conditions
on the ranges of swapped positions: a swap at any two positions preserves it.

\subsection{The trace equation system}

Symbolic execution of the outer loop (Definition 4.4 of the paper) introduces the synthetic counter
$K = \cnt(0, 1, n')$ of $n'$ iterations and one trace per modified variable. The inner loop is
summarized during the execution of the body (Theorem 4.8): its own trace equation system is in terms of
fresh variables, which become fresh \emph{sequences} over the outer iterations, and it is recorded as the
iteration constraint of the outer loop. After \texttt{-simplify} distributes the constraint over its
conjuncts, the goal is (variables renamed, the fragments of the inner loop grouped):
\begin{align}
  K &= \cnt(0, 1, n') \qquad I = \cnt(0, 1, n') \qquad \intf(n') = n \qquad b = \cur(B) \label{eq:outer}\\
  B &= \tr{a}{\bl{K}, \bl{I}, \seq{M}}{x, k, i, m \mapsto \swap{x}{k}{m}} \label{eq:B}\\[2pt]
  &\al{\bl{K}, \bl{I}, \seq{M}, \seq{M}'}{k, i, m, M \mapsto m = \cur(M)} \label{eq:f1}\\
  &\al{\bl{K}, \bl{I}, \bl{B}, \seq{M}', \seq{J}'}{k, i, x, M, J \mapsto {} \nonumber\\
  &\qquad\qquad M = \tr{k}{\bl{J}}{p, j \mapsto \upd{x[j] < x[p]}{j}\ \ \upd{x[p] \le x[j]}{p}}} \label{eq:f2}\\
  &\al{\bl{K}, \bl{I}, \seq{N}, \seq{K}'}{k, i, n'', K' \mapsto K' = \cnt(0, 1, n'')} \label{eq:f3}\\
  &\al{\bl{K}, \bl{I}, \seq{K}', \seq{J}'}{k, i, K', J \mapsto J = \sums(k + 1, \mapc{\bl{K'}}{1})} \label{eq:f4}\\
  &\al{\bl{K}, \bl{I}, \seq{J}', \seq{J}}{k, i, J, j \mapsto j = \cur(J)} \label{eq:f5}\\
  &\al{\bl{K}, \bl{I}, \seq{J}'}{k, i, J \mapsto \al{\bl{J}}{j \mapsto j < n}} \label{eq:f6}\\
  &\al{\bl{K}, \bl{I}, \seq{J}}{k, i, j \mapsto n \le j} \label{eq:f7}\\[2pt]
  &\Longrightarrow\ \sorted(b, 0, n) \wedge \perm(b, a) \nonumber
\end{align}
Equation \eqref{eq:B} is the outer array trace: in iteration $k$ (with $i = k$, since $I$ and $K$ are
equal counters that the simplifier does not identify), the positions $k$ and $m$ are swapped, where $m$ is
the result of the inner loop, the entry of the sequence $\seq{M}$ at $k$. The fragments
\eqref{eq:f1}--\eqref{eq:f7} describe that inner loop: its result $m$ is the current value of a
\emph{selection trace} $M$ \eqref{eq:f1}, which starts at $k$ and moves to $j$ whenever $x[j] < x[p]$,
over the inner counter $J$ \eqref{eq:f2}; $J$ is the trace of $j$, increased by 1 from $k + 1$ over the
synthetic inner counter $K'$ of $n''$ steps \eqref{eq:f3}, \eqref{eq:f4}, with the final value $j$
\eqref{eq:f5}; the test $j < n$ holds in every inner iteration \eqref{eq:f6} and fails at the exit
\eqref{eq:f7}. The generic fold of \texttt{teq.dfy} already turned the trace of $j$ into
$\sums(k + 1, \mapc{\bl{K'}}{1})$, since inside a fragment the equation $K' = \cnt(0, 1, n'')$ of the
sibling fragment \eqref{eq:f3} is not followed (problem~1 of nested loops in \texttt{NOTES.md}).

\subsection{Step 1: the inner counter (generic fragment lemmas)}\label{sec:sel:frag}

The two lemmas of \texttt{fragments.dfy} relate fragments of one inner loop to each other. They are
rewrite rules whose left-hand side is a fragment and whose side condition is a sibling fragment, matched
against the facts of the goal, which binds the sequences that the left-hand side does not mention.
\begin{align*}
  &\lemma{invs\_sums\_counter}:\\
  &\quad \al{ks, Ks, Ns}{k, K, n'' \mapsto K = \cnt(0, 1, n'')}\ \Longrightarrow \\
  &\al{ks, Ks, Js}{k, K, J \mapsto J = \sums(X(k), \mapc{\bl{K}}{c})}\\
  &\quad\Longleftrightarrow\ \al{ks, Ns, Js}{k, n'', J \mapsto J = \cnt(X(k), c, n'')}\\[4pt]
  &\lemma{invs\_current\_counter}:\\
  &\quad \al{ks, Ns, Js}{k, n'', J \mapsto J = \cnt(X(k), c, n'')}\ \Longrightarrow \\
  &\al{ks, Js, js}{k, J, j \mapsto j = \cur(J)}\\
  &\quad\Longleftrightarrow\ \al{ks, Ns, js}{k, n'', j \mapsto j = X(k) + \intf(n'') \cdot c}
\end{align*}
The initial value $X(k)$ is a higher-order pattern variable (Miller fragment): it matches $0$, $k$,
$k + 1$, or the value of a further dependency, so one lemma covers all inner counters. The first lemma
matches fragment \eqref{eq:f4} with $ks := \bl{K}$, $Ks := \seq{K}'$, $Js := \seq{J}'$,
$X := \lambda k.\, k + 1$, $c := 1$; its side condition is found as fragment \eqref{eq:f3}, which binds
$Ns := \seq{N}$. The unused dependency $\bl{I}$ of the fragment is selected by no pattern dependency and
dropped. The result replaces \eqref{eq:f4}:
\begin{equation}
  \al{\bl{K}, \seq{N}, \seq{J}'}{k, n'', J \mapsto J = \cnt(k + 1, 1, n'')}. \label{eq:f4b}
\end{equation}
The second lemma matches \eqref{eq:f5} with $Js := \seq{J}'$, $js := \seq{J}$, finds \eqref{eq:f4b} as its
side condition, and yields
\begin{equation}
  \al{\bl{K}, \seq{N}, \seq{J}}{k, n'', j \mapsto j = \intf(n'') + k + 1}. \label{eq:f5b}
\end{equation}
Conceptually, the fragment lemmas do by rewriting what the simplifier would do if the fragments of one
inner loop shared their facts (the sharing was implemented and reverted, since it made rewriting cycle):
they propagate the definition of a bound variable from one fragment into another. They are candidates for
\texttt{teq.dfy}; they close the two minimal nested-loop problems \texttt{invs\_counter\_equation.dfy} and
\texttt{invs\_fragments.dfy}, and \texttt{nested\_sum.dfy}, \texttt{nested\_mult.dfy},
\texttt{distinct.dfy} and the sorting algorithms use them.

\subsection{Step 2: the summary of the inner loop}\label{sec:sel:sum}

The inner loop finds the position of the minimum of $x[k..n)$. We name that position by an abstract
function $\argmin(x, lo, hi)$, declared without a definition, and state the summary as a rewrite rule
whose side conditions are the fragments, in the forms they have after Step~1:
\begin{align*}
  &\lemma{current\_selection\_swaps}:\\
  &\quad \al{ks, ms, Ms}{k, m, M \mapsto m = \cur(M)} \\
  &\quad\wedge\ \mathsf{all}(\mathsf{for}\ ks, \bl{\tr{a}{ks, is, ms}{x, k, i, m \mapsto \swap{x}{k}{m}}}, Ms, Js\\
  &\qquad\qquad \mathsf{check}\ k, x, M, J \mapsto {}\\
  &\qquad\qquad\qquad M = \tr{k}{\bl{J}}{p, j \mapsto \upd{x[j] < x[p]}{j}\ \ \upd{x[p] \le x[j]}{p}}) \\
  &\quad\wedge\ \al{ks, Ns, Js}{k, n'', J \mapsto J = \cnt(k + 1, 1, n'')} \\
  &\quad\wedge\ \al{ks, Ns, js}{k, n'', j \mapsto j = \intf(n'') + k + 1}\\
  &\quad\wedge\ \al{ks, js}{k, j \mapsto n \le j} \\
  &\quad\wedge\ \al{ks, Js_2}{k, J \mapsto \al{\bl{J}}{j \mapsto j < n_2}}\ \wedge\ Js_2 = Js\ \wedge\ n_2 = n \\
  &\quad\Longrightarrow\ \cur(\tr{a}{ks, is, ms}{x, k, i, m \mapsto \swap{x}{k}{m}})\\
  &\qquad\qquad = \cur(\tr{a}{ks}{x, k \mapsto \swap{x}{k}{\argmin(x, k, n)}})
\end{align*}
The left-hand side is matched against $b = \cur(B)$ with the definition \eqref{eq:B} of $B$, which binds
$a$, $ks := \bl{K}$, $is := \bl{I}$, $ms := \seq{M}$. The side conditions are then searched among the
facts in order, each binding a new sequence: \eqref{eq:f1} binds $Ms := \seq{M}'$; \eqref{eq:f2} binds
$Js := \seq{J}'$, where the dependency $\bl{\tr{a}{\ldots}}$ of the pattern is matched against $\bl{B}$
through the definition of $B$; \eqref{eq:f4b} binds $Ns := \seq{N}$; \eqref{eq:f5b} binds $js := \seq{J}$;
\eqref{eq:f7} binds $n$. The remaining condition is the inner test \eqref{eq:f6}, which is matched by a
duplicate $Js_2$ with a fresh bound $n_2$, identified afterwards by the equations $Js_2 = Js$ and $n_2 = n$:
a side condition whose sequence variables are all bound is \emph{decided} by the simplifier rather than
matched against the facts, and the decision fails, since the facts keep $\bl{I}$ as a dependency and
rename the bound variables. The equations are closed and trivially decided.

Soundness of the lemma: given the side conditions, in iteration $k$ the selection trace $M$ scans the
positions $k + 1, \ldots, k + n''$ with $n'' + k + 1 \ge n$ (exit) and $k + 1 + n'' - 1 < n$ if $n'' > 0$
(test), hence exactly $k + 1, \ldots, n - 1$, so $\cur(M)$ is a position of the minimum of $x[k..n)$,
i.e. $\argmin(x, k, n)$ for any choice of $\argmin$ among the minima. The result is the trace of swaps of
$k$ with $\argmin(x, k, n)$, the \emph{summary} of the whole algorithm as a trace of an abstract step.
After the rewrite, $b = \cur(\tr{a}{\bl{K}}{x, k \mapsto \swap{x}{k}{\argmin(x, k, n)}})$.

Two workarounds are visible: the inner test is kept in its raw form \eqref{eq:f6} as a side condition
instead of being resolved to a bound on $n''$, since the resolved form $k + 1 + \intf(n'') \le n$ is
matched by \lemma{above\_iterations\_def} of \texttt{teq.dfy} (pattern $k < t$ with the dependency $t$
synthesized as a map), which hides the fact; and the initial value $k + 1$ and the final value
$\intf(n'') + k + 1$ are written literally in the normal form of the simplifier, since a higher-order
summand $X(k)$ matches a single summand only.

\subsection{Step 3: sortedness by induction, permutation by a generic lemma}\label{sec:sel:ind}

The postcondition $\sorted(b, 0, n)$ is now
\[ \sorted(\cur(\tr{a}{\bl{K}}{x, k \mapsto \swap{x}{k}{\argmin(x, k, n)}}), 0, n). \]
The invariant of
selection sort after $k$ iterations is that the prefix $x[0..k)$ is sorted and all of it is at most all
of $x[k..n)$, written with a second abstract predicate $\partitioned(x, lo, mid, hi)$ (all of $x[lo..mid)$
are at most all of $x[mid..hi)$). The induction is a lemma over a trace of arbitrary steps $F$:
\begin{align*}
  &\lemma{sorted\_prefix\_steps}:\\
  &\quad  n = \intf(m)\ \wedge\
    \big(\forall x, k.\ 0 \le k < n \wedge \sorted(x, 0, k) \wedge \partitioned(x, 0, k, n) \\
  &\qquad\qquad \Longrightarrow \sorted(F(x, k), 0, k + 1) \wedge \partitioned(F(x, k), 0, k + 1, n)\big) \\
  &\quad\Longrightarrow\ \sorted(\cur(\tr{a}{\bl{\cnt(0, 1, m)}}{x, k \mapsto F(x, k)}), 0, n)\ \Longleftrightarrow\ \mathsf{true}
\end{align*}
The left-hand side matches with $F := \lambda x\, k.\ \swap{x}{k}{\argmin(x, k, n)}$ and $m := n'$,
through the definition $K = \cnt(0, 1, n')$ in the dependency position. The condition $n = \intf(m)$
is the fact $\intf(n') = n$ of \eqref{eq:outer}; it fixes the number of steps to exactly $n$, since
further steps would be unconstrained by the second condition (a lemma with $n \le \intf(m)$ would be
unsound). The quantified condition is decided by the simplifier: $F$ is beta-reduced, and the two
conjuncts of the conclusion are rules on atoms,
\begin{align*}
  &\lemma{sorted\_swap\_argmin}:\\
  &\quad k < n \wedge \sorted(x, 0, k) \wedge \partitioned(x, 0, k, n)\\
    &\quad\Longrightarrow \sorted(\swap{x}{k}{\argmin(x, k, n)}, 0, k + 1) \Longleftrightarrow \mathsf{true},\\
  &\lemma{partitioned\_swap\_argmin}:\\
  &\quad k < n \wedge \sorted(x, 0, k) \wedge \partitioned(x, 0, k, n)\\
    &\quad\Longrightarrow \partitioned(\swap{x}{k}{\argmin(x, k, n)}, 0, k + 1, n) \Longleftrightarrow \mathsf{true},
\end{align*}
whose side conditions are the hypotheses of the quantified implication. The base case is inside the
induction lemma: $\sorted(a, 0, 0)$ and $\partitioned(a, 0, 0, n)$ hold trivially (two further lemmas,
\lemma{sorted\_empty} and \lemma{partitioned\_empty}, state them for completeness). This is the method of
the paper's \lemma{current\_trace\_bounded}, an invariant scheme as a lemma with a quantified side
condition: the lemma database proves the induction step, the shape of the invariant is the lemma's
knowledge of selection sort, and no invariant inference takes place.

The postcondition $\perm(b, a)$ follows from a generic lemma on traces of swaps at positions that may
depend on the state and the iteration (a higher-order pattern):
\begin{align*}
  &\lemma{perm\_swaps\_one}:\\
  &\quad \perm(\cur(\tr{a}{ks}{x, k \mapsto \swap{x}{P(x, k)}{Q(x, k)}}), c)\\
  &\quad\Longleftrightarrow \perm(a, c)
\end{align*}
with $P := \lambda x\, k.\ k$ and $Q := \lambda x\, k.\ \argmin(x, k, n)$, and then \lemma{perm\_refl}.
The whole goal closes with \texttt{./Cuvee.sh -teq -simplify}; a wrong postcondition, e.g.
$\sorted(b, 0, n + 1)$, leaves a goal.

\subsection{What is generic and what is specific}

Of the lemmas used, \lemma{invs\_sums\_counter}, \lemma{invs\_current\_counter} and
\lemma{perm\_swaps\_one} are independent of the algorithm; \lemma{sorted\_prefix\_steps} is the
induction for every algorithm that grows a sorted prefix (selection sort, and insertion sort in a
variant with two dependencies); the summary \lemma{current\_selection\_swaps} and the two step lemmas are the
algorithm's knowledge: what one inner loop computes, and why one step keeps the invariant. The abstract
predicates and $\argmin$ are axiomatized by these lemmas only; a definition in Lean, as for
\texttt{teq.dfy}, would make the step lemmas and the summary provable.

\section{Bubble Sort: an Invariant Indexed by the Iteration}\label{sec:bs}

\subsection{The program}
\begin{lstlisting}
method main(a: array<int, int>, n: int) returns (b: array<int, int>)
  requires 0 <= n
  ensures sorted(b, 0, n)
  ensures perm(b, a)
{
  var i: int; var j: int; var t: int;
  b := a; i := n;
  while 1 < i {
    j := 0;
    while j + 1 < i {
      if b[j + 1] < b[j] { t := b[j]; b[j] := b[j + 1]; b[j + 1] := t; }
      j := j + 1;
    }
    i := i - 1;
  }
}
\end{lstlisting}

\subsection{The trace equation system}
The outer loop modifies only $b$, and only through the inner loop, so its trace is the chain of the inner
results, which the fold \lemma{results} names:
\begin{align}
  K &= \cnt(0, 1, n') \qquad I = \cnt(n, -1, n') \qquad b = \cur(B) \label{eq:bs0}\\
  B &= \mathsf{results}(a, \bl{K}, \seq{B}) \quad\text{where}\quad \mathsf{results}(x_0, ks, bs) = \tr{x_0}{ks, bs}{x, k, b \mapsto b} \nonumber\\
  &\al{\bl{K}, \bl{I}, \seq{B}, \seq{B}'}{k, i, b, B' \mapsto b = \cur(B')} \label{eq:bs1}\\
  &\al{\bl{K}, \bl{B}, \bl{I}, \seq{B}', \seq{J}'}{k, x, i, B', J \mapsto B' = \mathsf{bubble\_pass}(x, \bl{J})} \label{eq:bs2}\\
  &\qquad\text{where}\quad \mathsf{bubble\_pass}(b, js) = \tr{b}{js}{x, j \mapsto {}\nonumber\\
  &\qquad\qquad\qquad\qquad \upd{x[j+1] < x[j]}{\swap{x}{j}{j + 1}}\ \ \upd{x[j] \le x[j+1]}{x}} \nonumber\\
  &\al{\bl{K}, \bl{I}, \seq{J}}{k, i, j \mapsto n \le j + k + 1} \label{eq:bs3}\\
  &\al{\bl{K}, \bl{I}, \seq{J}'}{k, i, J \mapsto \al{\bl{J}}{j \mapsto j + k + 1 < n}} \label{eq:bs4}\\
  &\al{\bl{K}, \bl{I}, \seq{J}, \seq{J}'}{k, i, j, J \mapsto j = \cur(J)} \label{eq:bs5}\\
  &\al{\bl{K}, \bl{I}, \seq{K}', \seq{J}'}{k, i, K', J \mapsto J = \sums(0, \mapc{\bl{K'}}{1})} \label{eq:bs6}\\
  &\ldots \wedge\ n \le \intf(n') + 1\ \Longrightarrow\ \sorted(b, 0, n) \wedge \perm(b, a) \nonumber
\end{align}
The inner trace folds into $\mathsf{bubble\_pass}(x, \bl{J})$, a pass over the positions of the inner
counter $J$, by a \texttt{\{:fold\}} function of the file, like $\sums$ in \texttt{teq.dfy}. The test and
the exit are stated with $i = n - k$ already substituted: the pass of iteration $k$ runs over the
positions $0, \ldots, n - k - 2$.

\subsection{The summary}
\begin{align*}
  &\lemma{current\_results\_bubble}:\\
  &\quad \al{ks, bs, Bs}{k, b, B' \mapsto b = \cur(B')}\\
  &\quad\wedge\ \al{ks, \bl{\mathsf{results}(a, ks, bs)}, Bs, Js}{k, x, B', J \mapsto B' = \mathsf{bubble\_pass}(x, \bl{J})}\\
  &\quad\wedge\ \al{ks, js}{k, j \mapsto len \le j + k + 1}\\
  &\quad\wedge\ \al{ks, Js_2}{k, J \mapsto \al{\bl{J}}{j \mapsto j + k + 1 < len}}\ \wedge\ Js_2 = Js\\
  &\quad\wedge\ \al{ks, js_3, Js_3}{k, j, J \mapsto j = \cur(J)}\ \wedge\ js_3 = js\ \wedge\ Js_3 = Js\\
  &\quad\wedge\ \al{ks, Ks, Js}{k, K', J \mapsto J = \sums(0, \mapc{\bl{K'}}{1})}\\
  &\quad\Longrightarrow\ \cur(\mathsf{results}(a, ks, bs))\\
  &\qquad\qquad = \cur(\tr{a}{ks}{b, k \mapsto \cur(\mathsf{bubble\_pass}(b, \bl{\cnt(0, 1, \natf(len - k - 1))}))})
\end{align*}
This is the summary in its original form (October 9), before the fragment lemmas: the inner counter is
matched as $\sums$, and the test and exit bind $len$. The left-hand side $\cur(\mathsf{results}(a, ks, bs))$
is matched against $b = \cur(B)$ through \eqref{eq:bs0}, binding $ks := \bl{K}$ and $bs := \seq{B}$;
\eqref{eq:bs1} binds $Bs := \seq{B}'$, \eqref{eq:bs2} binds $Js := \seq{J}'$, \eqref{eq:bs3} binds
$js := \seq{J}$ and $len := n$; \eqref{eq:bs4}, \eqref{eq:bs5}, \eqref{eq:bs6} are matched through
duplicates and the equations identify them, as in Section~\ref{sec:sel:sum}. The result is the trace of passes, the
$k$-th over $\bl{\cnt(0, 1, \natf(n - k - 1))}$. Note that the arithmetic $n - k - 1$ on the bound
variable $k$ is admissible on the right-hand side of a rule; on a left-hand side it would not be matched,
which is why the original version of this file had a function $\mathsf{passlen}(len, k)$ for it, which
the induction below makes unnecessary.

\subsection{The induction with a clamped invariant}
After $k$ passes the suffix $b[n - k..n)$ is sorted and all of $b[0..n - k)$ are at most all of it. Stated
with $\sorted$ and $\partitioned$ as in Section~\ref{sec:sel:ind}, the invariant makes sense only for $k \le n$, but
the number of outer iterations $n'$ is known only as $n \le n' + 1$ together with the iteration
constraint $1 < i$, a disjunction $n' = 0 \vee (2 < n \wedge n' < n)$ that the simplifier does not split,
so the induction cannot fix the number of steps to $n - 1$ (as Section~\ref{sec:sel:ind} fixes it to $n$). Instead the
invariant is a predicate indexed by the iteration, $\mathsf{bubbled}(b, k, n)$, which means the sorted
suffix of length $k$ once $k \le n$ and the whole range sorted once $k \ge n$; then every step keeps it,
also the empty passes beyond $n - 1$:
\begin{align*}
  &\lemma{bubbled\_pass}:\\
  &\quad p = len - k - 1 \wedge \mathsf{bubbled}(b, k, len)\\
  &\quad\Longrightarrow \mathsf{bubbled}(\cur(\mathsf{bubble\_pass}(b, \bl{\cnt(0, 1, \natf(p))})), k + 1, len)
    \Longleftrightarrow \mathsf{true}\\[4pt]
  &\lemma{sorted\_bubbled\_steps}:\\
  &\quad len \le \intf(m) + 1\\
  &\quad\wedge\ \big(\forall b, k.\ 0 \le k < \intf(m) \wedge \mathsf{bubbled}(b, k, len) \Longrightarrow \mathsf{bubbled}(F(b, k), k + 1, len)\big)\\
  &\quad\Longrightarrow\ \sorted(\cur(\tr{a}{\bl{\cnt(0, 1, m)}}{b, k \mapsto F(b, k)}), 0, len)\ \Longleftrightarrow\ \mathsf{true}
\end{align*}
The induction matches with $F := \lambda b\, k.\ \cur(\mathsf{bubble\_pass}(b, \bl{\cnt(0, 1, \natf(n - k - 1))}))$,
$m := n'$, $len := n$. Its quantified condition is decided by \lemma{bubbled\_pass}, whose length $p$ is a
variable with the equation $p = len - k - 1$ as a side condition, since a pattern $len - k - 1$ with two
variables is not matched; the equation is closed after $p$ is bound to $n - k - 1$ by the left-hand side
and decided by arithmetic. The lemma is sound for any number of steps $m$ with $len \le m + 1$: the
invariant at $k = 0$ holds trivially, and $\mathsf{bubbled}(b, m, len)$ with $len - m \le 1$ is
$\sorted(b, 0, len)$. For $\perm$, \lemma{perm\_steps} (the quantified condition $\forall b, k.\ \perm(F(b, k), b)$,
decided by \lemma{perm\_bubble\_pass}) and \lemma{perm\_refl}.

\subsection{The idea}
Two things differ from selection sort. The inner loop is summarized to a function of the previous
array that is itself a trace, $\cur(\mathsf{bubble\_pass}(\ldots))$, so the outer trace is a trace of
traces; the lemma database needs no knowledge of what a pass does beyond the two facts
\lemma{bubbled\_pass} and \lemma{perm\_bubble\_pass}. And the invariant is written as one predicate over
the iteration, which absorbs the case distinction on the number of passes that the goal does not provide.
The pattern "invariant as a predicate of $(b, k)$, step lemma, generic induction" is the same as
\lemma{current\_trace\_bounded} of \texttt{teq.dfy}, with the induction hypothesis as an abstract atom.

\section{Cocktail Sort: Two Inner Loops and a Closed Side Condition}\label{sec:cock}

\subsection{The program}
\begin{lstlisting}
method main(a: array<int, int>, n: int) returns (b: array<int, int>)
  requires 0 <= n
  ensures sorted(b, 0, n)
  ensures perm(b, a)
{
  var lo: int;
  var hi: int;
  var j: int;
  var p: int;
  var t: int;
  b := a;
  lo := 0;
  hi := n;
  while lo + 1 < hi
  {
    j := lo;
    while j + 1 < hi
    {
      if b[j + 1] < b[j] {
        t := b[j];
        b[j] := b[j + 1];
        b[j + 1] := t;
      }
      j := j + 1;
    }
    hi := hi - 1;
    j := hi - 1;
    p := hi - 2;
    while lo < j
    {
      if b[j] < b[p] {
        t := b[j];
        b[j] := b[p];
        b[p] := t;
      }
      j := j - 1;
      p := p - 1;
    }
    lo := lo + 1;
  }
}
\end{lstlisting}
Each outer iteration runs a forward bubble pass over $[lo, hi)$, which moves the maximum to $hi - 1$, and a
backward pass over $(lo, hi - 1]$, which moves the minimum to $lo$; the backward pass uses two pointers
$j$ and $p = j - 1$, like \texttt{insertion\_sort2.dfy}, so that the positions of its swaps are counters.

\subsection{The trace equation system}
The outer loop modifies $b$ only through the two inner loops, so its trace is again the chain of the
results, $B = \mathsf{results}(a, \bl{K}, \seq{B}'')$, where the intermediate results $\seq{B}'$ of the forward
passes are the initial values of the backward passes. With $lo = k$ and $hi = n - k$ substituted,
\begin{align}
  K &= Lo = \cnt(0, 1, n') \qquad Hi = \cnt(n, -1, n') \qquad b = \cur(B) \qquad lo = \intf(n') \label{eq:ck0}\\
  &\al{\bl{K}, \bl{Lo}, \bl{Hi}}{k, lo, hi \mapsto 2k + 1 < n} \label{eq:ck1}\\
  &\al{\bl{K}, \bl{Lo}, \bl{Hi}, \seq{B}', \seq{B}'_t}{k, lo, hi, c, C \mapsto c = \cur(C)} \nonumber\\
  &\al{\bl{K}, \bl{B}, \bl{Lo}, \bl{Hi}, \seq{B}'_t, \seq{J}'_t}{k, x, lo, hi, C, J \mapsto C = \mathsf{bubble\_pass}(x, \bl{J})} \label{eq:ck2}\\
  &\al{\bl{K}, \seq{N}, \seq{J}'_t}{k, nn, J \mapsto J = \cnt(k, 1, nn)} \nonumber\\
  &\al{\bl{K}, \bl{Lo}, \bl{Hi}, \seq{J}'}{k, lo, hi, j \mapsto n \le j + k + 1} \qquad \ldots \label{eq:ck3}\\
  &\al{\bl{K}, \bl{Lo}, \bl{Hi}, \seq{B}'', \seq{B}''_t}{k, lo, hi, b, B' \mapsto b = \cur(B')} \nonumber\\
  &\al{\bl{K}, \bl{Lo}, \bl{Hi}, \seq{B}', \seq{B}''_t, \seq{J}''_t, \seq{P}''_t}{k, lo, hi, c, B', J, P \mapsto {}\nonumber\\
  &\qquad\qquad\qquad\qquad B' = \mathsf{bubble\_pass\_back}(c, \bl{J}, \bl{P})} \label{eq:ck4}\\
  &\al{\bl{K}, \seq{N}', \seq{J}''_t}{k, nn, J \mapsto J = \cnt(n - k - 2, -1, nn)} \nonumber\\
  &\al{\bl{K}, \seq{N}', \seq{P}''_t}{k, nn, P \mapsto P = \cnt(n - k - 3, -1, nn)} \label{eq:ck5}\\
  &\al{\bl{K}, \bl{Lo}, \bl{Hi}, \seq{J}''}{k, lo, hi, j \mapsto j \le k} \nonumber\\
  &\al{\bl{K}, \bl{Lo}, \bl{Hi}, \seq{J}'', \seq{J}''_t}{k, lo, hi, j, J \mapsto j = \cur(J)} \qquad \ldots \label{eq:ck6}\\
  &\ldots \wedge\ n \le \intf(n') + lo + 1\ \Longrightarrow\ \sorted(b, 0, n) \wedge \perm(b, a) \nonumber
\end{align}
with $\mathsf{bubble\_pass}$ as in the bubble sort and $\mathsf{bubble\_pass\_back}(c, js, ps)$ the trace of
conditional swaps $\swap{x}{j}{p}$ over the two counters. All facts of the outer iteration carry the
counters $Lo$ and $Hi$ as dependencies, which no body uses, and the outer exit is $n \le \intf(n') + lo + 1$
with $lo = \intf(n')$ a separate fact.

\subsection{The summary and the induction}
The summary \lemma{current\_results\_shake} takes the fragments of both inner loops as side conditions,
twenty-one of them, in the manner of Section~\ref{sec:sel:sum}: the results and traces of each pass, the
counters (whose initial values $n - k - 2$, $n - k - 3$ are higher-order $X(k)$, $Y(k)$, pinned by closed
$\mathsf{all}$ conditions $\al{ks}{k \mapsto X(k) = len - k - 2}$, since a sum of a rule variable and a bound
variable is not matched), the tests and the exits; the result is
\[ \cur(\mathsf{results}(a, ks, bs)) = \cur(\tr{a}{ks}{b, k \mapsto \mathsf{shake}(b, k, len)}) \]
with $\mathsf{shake}(b, k, len)$ abstract: the pass of iteration $k$. The induction is that of the bubble sort
with the invariant $\mathsf{shaken}(b, k, len)$, both ends of length $k$ sorted and the middle between them,
the whole range once $2k + 1 \ge len$, and \lemma{shaken\_shake} as the step.

The goal stayed open for a day although every piece was in place, and the diagnosis is the point of this
section. Truncated copies of the summary with its first $k$ side conditions showed that it fires with
$k \le 20$ and not with the 21st, $\al{ks, js_4, Js_3}{k, j, J \mapsto j = \cur(J)}$, the final values of
the backward counter. All its variables were bound by earlier conditions, so the condition was closed and
\emph{decided} by the simplifier instead of matched against the fact \eqref{eq:ck6}, and the simplifier
does not find a closed $\mathsf{all}$ among the facts (these carry the unused dependencies $\bl{Lo}$, $\bl{Hi}$,
and their bound variables are renamed). The same condition for the forward pass had fresh variables
$js_5, Js_5$ with identifying equations, the pattern of Section~\ref{sec:sel:sum}, and succeeded. With
fresh $js_6, Js_6$ and $js_6 = js_4 \wedge Js_6 = Js_3$ the summary fires and $\perm$ closes.

Then $\sorted(b, 0, n)$ remained: the induction's bound $len \le 2\,\intf(m) + 1$ is not a difference
constraint, and the simplifier does not derive it from the two facts $n \le \intf(n') + lo + 1$ and
$lo = \intf(n')$. The lemma therefore states the bound in the goal's form,
\begin{align*}
  &\lemma{sorted\_shaken\_steps}:\quad l = \intf(m)\ \wedge\ len \le \intf(m) + l + 1\\
  &\quad\wedge\ \big(\forall b, k.\ 0 \le k < \intf(m) \wedge \mathsf{shaken}(b, k, len) \Longrightarrow \mathsf{shaken}(F(b, k), k + 1, len)\big)\\
  &\quad\Longrightarrow\ \sorted(\cur(\tr{a}{\bl{\cnt(0, 1, m)}}{b, k \mapsto F(b, k)}), 0, len)\ \Longleftrightarrow\ \mathsf{true}
\end{align*}
with $l$ an open variable matched against the equation of $lo$, after which the bound is literally a fact.
The counter $\bl{\cnt(0, 1, m)}$ as a dependency matches $\bl{K}$ through $K = \cnt(0, 1, n')$ in the
context, as in the bubble sort. Two things that were suspected were not the cause: the extra dependencies
$\bl{Lo}$, $\bl{Hi}$ of the facts (a diagnostic with two and with four dependencies matched both) and the
name $n$ of the rule variable, which the program shares (it was renamed to $len$ anyway). The file in the
tree now runs with the module of the fragments of inner loops (\texttt{-include teq:nested}), which resolves
the inner counters to $\cnt(k, 1, nn)$ and their final values to $j = \intf(nn) + k$,
$\intf(nn) + j + k + 2 = len$ and $\intf(nn) + j + k + 3 = len$; the summary states these conditions in the
forms the module produces, as the goal prints them, instead of $j = \cur(J)$. The two insertion sorts were
adapted the same way, so that every nested loop of the folder verifies with the module.

\subsection{The idea}
Nothing new is needed for two inner loops: each pass is summarized by its own fold, and the summary of
the outer iteration takes both. The lesson is about side conditions: a condition whose variables are all
bound is decided, not matched, so every $\mathsf{all}$ condition that should be read off a fact must bind a
fresh sequence and be identified afterwards; and a closed arithmetic condition must be the fact as the goal
has it, since the decision procedure combines difference constraints only.

\section{Merge: Pinned Guards and Open Trace Variables}\label{sec:merge}

\subsection{The program}
\begin{lstlisting}
method main(a: array<int, int>, m: int, b: array<int, int>, n: int) returns (c: array<int, int>)
  requires 0 <= m && 0 <= n
  requires sorted(a, 0, m) && sorted(b, 0, n)
  ensures sorted(c, 0, m + n)
  ensures forall x: int :: count(x, selects(c, butlast(counter(0, 1, natural(m + n)))))
    == count(x, selects(a, butlast(counter(0, 1, natural(m)))))
     + count(x, selects(b, butlast(counter(0, 1, natural(n)))))
{
  var i: int; var j: int; var k: int;
  i := 0; j := 0; k := 0;
  while k < m + n {
    if n <= j || (i < m && a[i] <= b[j]) { c := c[k := a[i]]; i := i + 1; }
    else { c := c[k := b[j]]; j := j + 1; }
    k := k + 1;
  }
}
\end{lstlisting}
The second postcondition states the permutation as a multiset: for each value, its number of occurrences
among the first $m + n$ elements of $c$ is the sum of its occurrences in $a$ and $b$, in the vocabulary of
\texttt{teq.dfy}.

\subsection{The trace equation system}
\begin{align}
  K &= \cnt(0, 1, n') \qquad c = \cur(C) \qquad k = \intf(n') \qquad m + n \le k \label{eq:mg0}\\
  C &= \tr{c_0}{\bl{K}, \bl{I}, \bl{J}}{c, k, i, j \mapsto \upd{G}{c[k := a[i]]}\ \ \upd{\neg G}{c[k := b[j]]}} \label{eq:mg1}\\
  I &= \tr{0}{\bl{K}, \bl{J}}{i, k, j \mapsto \upd{G}{i + 1}\ \ \upd{\neg G}{i}} \label{eq:mg2}\\
  J &= \tr{0}{\bl{K}, \bl{I}}{j, k, i \mapsto \upd{G}{j}\ \ \upd{\neg G}{j + 1}} \label{eq:mg3}
\end{align}
where $G$ abbreviates the test $n \le j \vee (i < m \wedge a[i] \le b[j])$ of the program. The pointers
are traces over each other, \eqref{eq:mg2} and \eqref{eq:mg3}: $I$ has $\bl{J}$ among its dependencies and $J$ has $\bl{I}$. A single loop,
no fragments; the whole difficulty is in matching this shape.

\subsection{The lemmas}\label{sec:merge:lemmas}
The merge is named by an abstract list $\mathsf{merged}(a, m, b, n)$, the values of the result on
$0, \ldots, m + n - 1$, with two facts: it is sorted, and its counts are those of both ranges. The
lemma on the loop rewrites the values of the result:
\begin{align*}
  &\lemma{selects\_merge\_loop}:\\
  &\quad s = m + n\ \wedge\ m + n \le \intf(nn)\ \wedge\ 0 \le m \wedge 0 \le n\\
  &\quad\wedge\ (\forall i, j.\ G(i, j) \Leftrightarrow (n \le j \vee (i < m \wedge a[i] \le b[j])))\\
  &\quad\wedge\ (\forall i, j.\ U(i, j) = a[i])\ \wedge\ (\forall i, j.\ V(i, j) = b[j])\\
  &\quad\wedge\ Iv = \tr{0}{\bl{\cnt(0, 1, nn)}, js}{i, k, j \mapsto \upd{G(i, j)}{i + 1}\ \ \upd{\neg G(i, j)}{i}}\\
  &\quad\wedge\ is = \bl{Iv}\\
  &\quad\wedge\ Jv = \tr{0}{\bl{\cnt(0, 1, nn)}, is}{j, k, i \mapsto \upd{G(i, j)}{j}\ \ \upd{\neg G(i, j)}{j + 1}}\\
  &\quad\wedge\ js = \bl{Jv}\\
  &\quad\Longrightarrow\ \mathsf{selects}(\cur(\tr{c_0}{\bl{\cnt(0, 1, nn)}, is, js}{c, k, i, j \mapsto {}\\
  &\qquad\qquad\qquad \upd{G(i, j)}{c[k := U(i, j)]}\ \ \upd{\neg G(i, j)}{c[k := V(i, j)]}}), \bl{\cnt(0, 1, \natf(s))})\\
  &\qquad\qquad = \mathsf{merged}(a, m, b, n)
\end{align*}
and \lemma{sorted\_merge\_loop} is the same with the conclusion $\sorted(\cur(\ldots), 0, s) \Leftrightarrow \mathsf{true}$
and the two $\sorted$ premises; \lemma{count\_merged} gives
\begin{align*}
  \mathsf{count}(x, \mathsf{merged}(a, m, b, n)) = {}&\mathsf{count}(x, \mathsf{selects}(a, \bl{\cnt(0, 1, \natf(m))}))\\
  &+ \mathsf{count}(x, \mathsf{selects}(b, \bl{\cnt(0, 1, \natf(n))})).
\end{align*}

Three workarounds of the matcher are built into this lemma, each isolated by experiments
(\texttt{NOTES.md}, findings 14, 16 and the section on the pin):
\begin{enumerate}[nosep]
\item \textbf{The pin.} A literal disjunctive guard $n \le j \vee (\ldots)$ in a pattern is never
  matched. The guard is the higher-order variable $G(i, j)$, which matches any guard over the bound
  variables $i, j$, and the quantified side condition
  $\forall i, j.\ G(i, j) \Leftrightarrow (n \le j \vee \ldots)$ pins it to the intended one; the simplifier
  decides the pin by beta reduction. The stored values $U(i, j) = a[i]$ and $V(i, j) = b[j]$ are pinned
  likewise. (A higher-order variable for an array read $a[i]$ alone would be instantiated by the array
  itself, finding 10, hence the pin instead of a literal $a[i]$.)
\item \textbf{The open trace variable.} The lemma must know that $is$ and $js$ are the pointer traces,
  which the pattern cannot say, since each is defined over the other. A closed side condition
  $is = \bl{\tr{0}{\ldots}{\ldots}}$ is decided by the simplifier and fails (it does not expand $\bl{I}$).
  Instead $Iv = \tr{0}{\ldots, js}{\ldots}$ with an \emph{open} trace variable $Iv$ is matched against
  the hypothesis $I = \tr{0}{\bl{K}, \bl{J}}{\ldots}$ \eqref{eq:mg2}, which binds $Iv := I$ (and the
  counter $\bl{\cnt(0, 1, nn)}$ matches $\bl{K}$ through its definition); then $is = \bl{Iv}$ is closed
  and trivially true.
\item \textbf{The whole atom as the left-hand side.} A variable argument of an uninterpreted function
  is not expanded by its definition when the simplifier rewrites below it: $\sorted(c, 0, m + n)$ with
  $c = \cur(C)$ is not matched by a left-hand side $\cur(\tr{\ldots}{\ldots})$, but
  $\sorted(\cur(\tr{\ldots}{\ldots}), 0, s)$ as the whole atom is, and $\mathsf{selects}(\cur(\ldots), \ldots)$
  with $\mathsf{selects}$ as the head, inside $\mathsf{count}$ and the quantifier. The length is a
  variable $s$ with $s = m + n$ as a side condition, since a pattern $m + n$ of two variables is not matched.
\end{enumerate}
The remaining condition $m + n \le \intf(nn)$ is the exit $m + n \le k$ with $k = \intf(n')$ of
\eqref{eq:mg0}: after at least $m + n$ iterations the positions $0, \ldots, m + n - 1$ are those of the
merge, whatever further iterations store beyond them, which makes the lemma sound without fixing the
number of iterations exactly (the goal knows it only as a disjunction).

\subsection{The idea}
The loop computes a function of its inputs that has a name in the specification's vocabulary,
$\mathsf{merged}$, and the lemma database relates the loop to that name and the name to the properties.
The matching problems of mutually dependent traces are solved by two general devices, the pin and the
open trace variable, which apply to every literal pattern that the matcher does not handle: in
\texttt{insertion\_sort.dfy} with one pointer, the positions $j$ and $j - 1$ of the swaps are pinned in
the same way, which closed a file that had been open since the start.

\section{Gnome Sort: a Data-Dependent Number of Iterations}\label{sec:gnome}

\subsection{The program}
\begin{lstlisting}
method main(a: array<int, int>, n: int) returns (b: array<int, int>)
  requires 0 <= n
  ensures sorted(b, 0, n)
  ensures perm(b, a)
{
  var i: int;
  var t: int;
  b := a;
  i := 1;
  while i < n
  {
    if i == 0 || b[i - 1] <= b[i] {
      i := i + 1;
    } else {
      t := b[i];
      b[i] := b[i - 1];
      b[i - 1] := t;
      i := i - 1;
    }
  }
}
\end{lstlisting}
The position moves up when the adjacent elements are in order and down with a swap otherwise. The number of
iterations is not a counter and is unrelated to $n$: the exit $n \le i$ bounds it from below only. The file
was the folder's negative example, "open, as expected", until it was verified with two lemmas.

\subsection{The trace equation system}
\begin{align}
  K &= \cnt(0, 1, n') \qquad i = \cur(I) \qquad b = \cur(B) \qquad n \le i \label{eq:gn0}\\
  I &= \tr{1}{\bl{K}, \bl{B}}{i, k, b \mapsto \upd{G}{i + 1}\ \ \upd{\neg G}{i - 1}} \label{eq:gn1}\\
  B &= \tr{a}{\bl{K}, \bl{I}}{b, k, i \mapsto \upd{G}{b}\ \ \upd{\neg G}{b[i := b[i - 1]][i - 1 := b[i]]}} \label{eq:gn2}\\
  &\qquad\text{where } G \text{ abbreviates } i = 0 \vee b[i - 1] \le b[i] \nonumber\\
  &\al{\bl{K}, \bl{I}}{k, i \mapsto i < n} \label{eq:gn3}
\end{align}
As in the merge and the lower bound, two traces over each other with a data-dependent guard.

\subsection{The lemmas}
The invariant of gnome sort is $\sorted(b, 0, i)$: a step up extends the sorted prefix by the element that
is in order with its predecessor, a swap down keeps the prefix below the position sorted, and at the exit
$n \le i$ the prefix covers the range. The lemma states this on the whole atom of the postcondition
(Section~\ref{sec:merge}), with the guard and the swap positions higher-order and pinned, since the literal
$i - 1$ is not matched:
\begin{align*}
  &\lemma{sorted\_gnome}:\quad Bv = \tr{a}{\bl{Kv}, \bl{Iv_2}}{b, k, i \mapsto \upd{G(i, b)}{b}\ \ \upd{\neg G(i, b)}{b[Q(i) := b[P(i)]][P(i) := b[Q(i)]]}}\\
  &\quad\wedge\ Iv_2 = Iv\\
  &\quad\wedge\ Iv_2 = \tr{1}{\bl{Kv_2}, \bl{Bv_2}}{i, k, b \mapsto \upd{G(i, b)}{i + u}\ \ \upd{\neg G(i, b)}{i + d}}\\
  &\quad\wedge\ Kv_2 = Kv \wedge Bv_2 = Bv \wedge u = 1 \wedge d = -1\\
  &\quad\wedge\ (\forall i, b.\ G(i, b) \Leftrightarrow (i = 0 \vee b[i - 1] \le b[i]))\ \wedge\ (\forall i.\ P(i) = i - 1)\ \wedge\ (\forall i.\ Q(i) = i)\\
  &\quad\wedge\ x = \cur(Iv)\ \wedge\ n \le x\\
  &\quad\Longrightarrow\ \sorted(\cur(\tr{a}{\bl{Kv}, \bl{Iv}}{b, k, i \mapsto \upd{G(i, b)}{b}\ \ \upd{\neg G(i, b)}{b[Q(i) := b[P(i)]][P(i) := b[Q(i)]]}}), 0, n)\
    \Longleftrightarrow\ \mathsf{true}
\end{align*}
The left-hand side matches $\sorted(b, 0, n)$ through $b = \cur(B)$ and \eqref{eq:gn2}, binding $Kv := K$,
$Iv := I$, $G$, $P := \lambda i.\ i - 1$, $Q := \lambda i.\ i$ by abstraction. The conditions follow the
discipline of Section~\ref{sec:fd} and the lower bound: the trace of the array is tied first, open in $Bv$
and $Iv_2$, matched against \eqref{eq:gn2}, and $Iv_2 = Iv$ is decided; the trace of the position is tied
second, with fresh $Kv_2$, $Bv_2$ and the steps $u$, $d$ as variables (a literal $-1$ in a pattern is not
matched), matched against \eqref{eq:gn1}, and identified; the pins are decided by beta reduction; and the
exit is read from the facts through $x = \cur(Iv)$, matched against $i = \cur(I)$, then $n \le x$ decided.
No condition mentions the left-hand side, so nothing re-enters. For the permutation, \lemma{perm\_gnome}
states that every step of such a trace is the identity or a swap, for any guard and positions, hence
$\perm(\cur(\ldots), c) \Leftrightarrow \perm(a, c)$, and \lemma{perm\_refl} closes it. A wrong postcondition
$\sorted(b, 0, n + 1)$ leaves the goal.

\subsection{The idea}
A data-dependent number of iterations is no obstacle when the property can be stated on the final values
alone: the lemma does not describe how many steps were taken, only that the invariant holds after any number
of them and that the exit turns it into the postcondition. This is the same shape as the merge, where the
pointers are traces over each other, and it shows that the "negative example" of the paper's loop shapes is
negative for the characterization of the traces, not for the verification of the postcondition.

\section{Counting Sort: Loops in Sequence and Atoms for Hypotheses}\label{sec:cs}

\subsection{The program}
\begin{lstlisting}
method main(a: array<int, int>, n: int, k: int) returns (b: array<int, int>)
  requires 0 <= n && 0 < k
  requires keys_below(a, n, k)              // forall i :: 0 <= i < n ==> 0 <= a[i] < k
  ensures sorted(b, 0, n)
  ensures forall x: int :: count(x, selects(b, butlast(counter(0, 1, natural(n)))))
    == count(x, selects(a, butlast(counter(0, 1, natural(n)))))
{
  var c: array<int, int>; var i: int; var v: int; var s: int; var t: int;
  v := 0;
  while v < k { c := c[v := 0]; v := v + 1; }
  i := 0;
  while i < n { c := c[a[i] := c[a[i]] + 1]; i := i + 1; }
  s := 0; v := 0;
  while v < k { t := c[v]; c := c[v := s]; s := s + t; v := v + 1; }
  i := 0;
  while i < n { b := b[c[a[i]] := a[i]]; c := c[a[i] := c[a[i]] + 1]; i := i + 1; }
}
\end{lstlisting}
Four loops in sequence: the histogram is zeroed, filled, replaced by its prefix sums (the start position
of each value), and the elements are scattered to their start positions, which are incremented.

\subsection{The trace equation system}
Each loop contributes a trace whose current value is the initial value of the next; the simplifier has
already applied the introduction rules of \texttt{teq.dfy} (stores, consts, sums, selects) when the goal
is printed ($V, V_2, K_2$ are counters of $k'$ steps with $\intf(k') = k$, $K, I, K_3, I_3$ counters of
$n'$ steps with $\intf(n') = n$):
\begin{align}
  c_1 &= \cur(\mathsf{stores}(c_0, \bl{V}, \bl{\mathsf{consts}(k', 0)})) \label{eq:cs1}\\
  c_2 &= \cur(\tr{c_1}{\bl{K}, \bl{I}}{c, k, i \mapsto c[a[k] := c[a[k]] + 1]}) \label{eq:cs2}\\
  c_3 &= \cur(\mathsf{stores}(c_2, \bl{V_2}, \bl{S})) \qquad S = \sums(0, \mathsf{selects}(c_2, \bl{K_2})) \label{eq:cs3}\\
  b' &= \cur(\mathsf{stores}(b, \mapc{\bl{K_3}, \bl{C}, \bl{I_3}}{k, c, i \mapsto c[a[k]]}, \mathsf{selects}(a, \bl{I_3}))) \label{eq:cs4}\\
  C &= \tr{c_3}{\bl{K_3}, \bl{I_3}}{c, k, i \mapsto c[a[k] := c[a[k]] + 1]} \label{eq:cs5}
\end{align}
The scatter \eqref{eq:cs4} is a $\mathsf{stores}$ whose keys are the positions read from the table $C$ as
it is incremented, and whose values are the elements of $a$.

\subsection{The lemma chain}
The histogram is characterized at each value (the generic lemma \lemma{current\_trace\_add\_at} of
\texttt{teq.dfy} is for a fixed key; here the key is $a[k]$):
\begin{align*}
  &\lemma{current\_trace\_histogram}:\\
  &\quad \cur(\tr{c_0}{\bl{\cnt(0, 1, n)}, hs}{c, k, h \mapsto c[a[k] := c[a[k]] + 1]})[v]\\
  &\qquad = c_0[v] + \mathsf{count}(v, \mathsf{selects}(a, \bl{\cnt(0, 1, n)}))
\end{align*}
It fires on the reads $c_2[v]$ inside $\mathsf{selects}(c_2, \bl{K_2})$ of \eqref{eq:cs3}, with
$c_1[v] = 0$ from \eqref{eq:cs1} by \lemma{current\_stores} of \texttt{teq.dfy}, so the prefix sums
become
\[ S = \sums(0, \mapc{\bl{K_2}, \bl{V_2}}{v, h \mapsto \mathsf{count}(v, \mathsf{selects}(a, \bl{\cnt(0, 1, n')}))}). \]
The table of start positions is then named by a lemma whose left-hand side is this whole term:
\begin{align*}
  &\lemma{rank\_table\_def}:\\
  &\quad \cur(\mathsf{stores}(c_0, \bl{\cnt(0, 1, kk)}, \bl{\sums(0, \mapc{\bl{\cnt(0, 1, kk)}, hs}{v, h \mapsto {}\\
  &\qquad\qquad\qquad \mathsf{count}(v, \mathsf{selects}(a, \bl{\cnt(0, 1, n)}))})}))
    = \mathsf{rank\_table}(c_0, a, n, kk)
\end{align*}
which rewrites $c_3$ to $\mathsf{rank\_table}(c_2, a, n', k')$, the number of elements of $a[0..n)$ below
each value $v < k$ (and $c_2$ beyond). Naming the table is essential: the scatter lemma below needs the
bound $kk$ of the keys, and a side condition $\mathsf{ranks}(c_0, a, n, kk)$ with $kk$ open would be
matched by no fact, while $\mathsf{rank\_table}(c_2, a, n, kk)$ in the pattern binds $kk$. Finally the
scatter:
\begin{align*}
  &\lemma{sorted\_scatter}:\\
  &\quad s = \intf(n)\ \wedge\ \mathsf{keys\_below}(a, s, kb)\ \wedge\ kb = \intf(kk)
    \ \wedge\ Iv = \cnt(0, 1, n)\ \wedge\ hs_3 = \bl{Iv}\\
  &\quad\Longrightarrow\ \sorted(\cur(\mathsf{stores}(b_0,
    \mapc{\bl{\cnt(0, 1, n)}, \bl{T}, hs_2}{k, c, h \mapsto c[a[k]]},\\
  &\qquad\qquad\qquad\qquad\qquad \mathsf{selects}(a, hs_3))), 0, s)\ \Longleftrightarrow\ \mathsf{true}\\
  &\qquad \text{where } T = \mathsf{trace}(\mathsf{rank\_table}(c_2, a, n, kk),\\
  &\qquad\qquad\qquad \mathsf{for}\ \bl{\cnt(0, 1, n)}, hs\ \mathsf{do}\ c, k, h \mapsto c[a[k] := c[a[k]] + 1])
\end{align*}
and \lemma{count\_scatter} with the conclusion
\[ \mathsf{count}(x, \mathsf{selects}(\cur(\ldots), \bl{\cnt(0, 1, n)})) = \mathsf{count}(x, \mathsf{selects}(a, \bl{\cnt(0, 1, n)})). \]
The left-hand side binds $n := n'$, $kk := k'$, $hs_3 := \bl{I_3}$ and the rest; the counter of the
values' $\mathsf{selects}$ is tied by the open trace variable $Iv$ (Section~\ref{sec:merge:lemmas}), since $\bl{I_3}$ and
$\bl{\cnt(0, 1, n')}$ are not identified by a closed equation. The precondition on the keys is the atom
$\mathsf{keys\_below}$ of the specification with an open bound $kb$, matched against the fact and then
identified with $\intf(kk)$: a quantified hypothesis $\forall i.\ 0 \le i < n \Rightarrow 0 \le a[i] < k$
is matched by no side condition, open or closed, so the specification uses the atom (as
\texttt{count\_occurrences.dfy} uses the vocabulary of \texttt{teq.dfy} in its postcondition).

Soundness of \lemma{sorted\_scatter}: with the start positions correct for all keys, element $a[k]$ is
stored at (number of elements below $a[k]$) + (number of equal elements before position $k$), which is
a bijection of $0, \ldots, n - 1$ that is monotone in the value, so the result is sorted on $0, \ldots, n - 1$
and has the elements of $a$. Both are asserted, not proved; the lemma is the algorithm's correctness
argument in one statement, over the exact shape of the two last loops.

\subsection{The idea}
Loops in sequence compose through the current values: each loop's trace has the previous result as its
initial value, and the acceleration lemmas of \texttt{teq.dfy} (stores read at a position, sums, consts)
push the characterization of one loop into the next without any lemma specific to counting sort, up to
the table of start positions. The algorithm-specific knowledge is concentrated in naming the table and
in the scatter lemma, whose left-hand side is the composition of the last two loops. What could not be
expressed is a quantified precondition as a side condition; the specification's vocabulary must provide
an atom for it.

\section{First Duplicate: a Search with an Early Exit}\label{sec:fd}

\subsection{The program}
\begin{lstlisting}
method main(a: array<int, int>, n: int) returns (r: bool)
  requires 0 <= n
  ensures r ==> !distincts(selects(a, butlast(counter(0, 1, natural(n)))))
  ensures !r ==> distincts(selects(a, butlast(counter(0, 1, natural(n)))))
{
  var i: int; var j: int;
  r := false; i := 0;
  while i < n && !r {
    j := i + 1;
    while j < n && !r {
      if a[i] == a[j] { r := true; }
      j := j + 1;
    }
    i := i + 1;
  }
}
\end{lstlisting}
The nested loop of \texttt{distinct.dfy} (which is verified) with both loops stopped by the flag: the
search for a duplicate ends at the first one. The postcondition uses $\mathsf{distincts}$ of the lemma
database, and is stated as two implications rather than the equivalence $r \Leftrightarrow \neg\mathsf{distincts}(\ldots)$,
for a reason explained below.

\subsection{The trace equation system}
The outer loop's constraint includes the flag, $\al{\bl{K}, \bl{R}, \bl{I}}{k, r, i \mapsto \neg r}$, so
$r$ is false at the head of every outer iteration, and the simplifier uses this inside the inner trace:
\begin{align}
  K &= I = \cnt(0, 1, n') \qquad r = \cur(R) \qquad R = \mathsf{results}(\mathsf{false}, \bl{K}, \seq{R}) \qquad i = \intf(n') \label{eq:fd0}\\
  &\al{\bl{K}, \bl{R}, \bl{I}}{k, r, i \mapsto \neg r} \label{eq:fd1}\\
  &\al{\bl{K}, \bl{I}, \seq{R}, \seq{R}'}{k, i, r, R' \mapsto r = \cur(R')} \label{eq:fd2}\\
  &\al{\bl{K}, \bl{I}, \seq{R}', \seq{J}'}{k, i, R', J \mapsto {}\nonumber\\
  &\qquad\qquad R' = \tr{\mathsf{false}}{\bl{J}}{r, j \mapsto \upd{a[k] = a[j]}{\mathsf{true}}\ \ \upd{a[k] \ne a[j]}{\mathsf{false}}}} \label{eq:fd3}\\
  &\al{\bl{K}, \seq{N}, \seq{J}'}{k, n'', J \mapsto J = \cnt(k + 1, 1, n'')} \label{eq:fd4}\\
  &\al{\bl{K}, \seq{N}, \seq{J}}{k, n'', j \mapsto j = \intf(n'') + k + 1} \label{eq:fd4b}\\
  &\al{\bl{K}, \bl{I}, \seq{J}'}{k, i, J \mapsto \al{\bl{J}}{j \mapsto j < n}} \label{eq:fd5}\\
  &\al{\bl{K}, \bl{I}, \seq{R}'}{k, i, R' \mapsto \al{\bl{R'}}{r \mapsto \neg r}} \label{eq:fd5b}\\
  &\al{\bl{K}, \bl{I}, \seq{R}, \seq{J}}{k, i, r, j \mapsto \neg(j < n \wedge \neg r)} \label{eq:fd6}\\
  &\neg(i < n \wedge \neg r)\ \Longrightarrow\ r \Leftrightarrow \neg\mathsf{distincts}(\mathsf{selects}(a, \bl{\cnt(0, 1, \natf(n))})) \nonumber
\end{align}
(after the fragment lemmas of Section~\ref{sec:sel:frag}, which gave \eqref{eq:fd4} and \eqref{eq:fd4b}). The inner update
\eqref{eq:fd3} does not mention the previous value of the flag, since the inner constraint
$\neg r$ \eqref{eq:fd5b} replaced it by $\mathsf{false}$: the inner trace is a sequence of comparisons
with $a[k]$. The inner exit \eqref{eq:fd6} is the negation of the test $j < n \wedge \neg r$: either all
later positions were visited, or the flag was set. The outer exit $\neg(i < n \wedge \neg r)$ is a
disjunction, $n \le i \vee r$, with one goal per postcondition.

\subsection{The lemmas}
The summary of the inner loop in the style of Section~\ref{sec:sel:sum} is a rewrite of the fragment \eqref{eq:fd2}
into a fact about the results: with \eqref{eq:fd3}--\eqref{eq:fd6} as
side conditions (duplicated and identified by equations where closed),
\begin{align*}
  \al{ks, rs, Rs}{k, r, R' \mapsto r = \cur(R')}\ \Longleftrightarrow\ \al{ks, rs}{k, r \mapsto r = \neg\mathsf{nodup}(a, k, n)}
\end{align*}
where $\mathsf{nodup}(a, k, n)$ is the abstract "$a[k]$ does not occur among $a[k + 1..n)$". This is
sound in both exit cases of the inner loop: if the flag was set, a later duplicate was found; if $j$
reached $n$, all later positions were compared. The outer trace is then
$\mathsf{results}(\mathsf{false}, \bl{K}, \seq{R})$ with $\seq{R}$ characterized pointwise, and the
simplifier itself turns it into $\mathsf{results}(\mathsf{false}, \bl{K}, \mapc{\bl{K}}{k \mapsto \neg\mathsf{nodup}(a, k, n)})$.

The outer loop is summarized on the atom $\mathsf{distincts}(\ldots)$ of the postcondition, by one rule per
case of the exit. If the flag is set, it is the result of the last iteration, which found a duplicate:
\begin{align*}
  &\lemma{distincts\_flag\_set}:\quad Kv = \cnt(0, 1, m)\\
  &\quad\wedge\ flag = \cur(\mathsf{results}(\mathsf{false}, \bl{Kv}, \mapc{\bl{Kv}}{k \mapsto \neg\mathsf{nodup}(a, k, n)}))\\
  &\quad\wedge\ flag\ \wedge\ \intf(m) \le n\ \Longrightarrow\ \mathsf{distincts}(\mathsf{selects}(a, \bl{\cnt(0, 1, \natf(n))}))\ \Longleftrightarrow\ \mathsf{false}
\end{align*}
If all iterations ran, $n \le \intf(m)$, and the flag was reset at the head of each, \eqref{eq:fd1}, then no
element before the last has a later duplicate (the flag at the head of iteration $k$ is the result of
iteration $k - 1$) and the last has none:
\begin{align*}
  &\lemma{distincts\_flags\_reset}:\quad Kv = \cnt(0, 1, m)\ \wedge\ Rv = \mathsf{results}(\mathsf{false}, \bl{Kv}, rs)\\
  &\quad\wedge\ \al{\bl{Kv}, hs}{k, h \mapsto \neg h}\ \wedge\ hs = \bl{Rv}\\
  &\quad\wedge\ \al{\bl{Kv}, rs_2}{k, h \mapsto h = \neg\mathsf{nodup}(a, k, n)}\ \wedge\ rs_2 = rs\ \wedge\ n \le \intf(m)\\
  &\quad\Longrightarrow\ \mathsf{distincts}(\mathsf{selects}(a, \bl{\cnt(0, 1, \natf(n))}))\ \Longleftrightarrow\ \mathsf{true}
\end{align*}
Three details made these fire. The outer counter and the flag trace are the variables $K$ and $R$ of the
goal with their equations \eqref{eq:fd0} as hypotheses, so a rule with $\bl{\cnt(0, 1, m)}$ as an
\emph{argument} never matches; the open trace variables $Kv$, $Rv$ are matched against the hypotheses
(Section~\ref{sec:merge}). The fact of the flags at the heads \eqref{eq:fd1} carries the counter $\bl{I}$ as
an unused dependency, which the simplifier drops from a closed instance of the condition, after which it
is not found among the facts; hence the open list $hs$, identified with $\bl{Rv}$ afterwards. And the exit
$n \le i \vee r$ is needed as a closed condition: with the postcondition as an equivalence, the goal's
conclusion is $i < n \wedge \neg r\ \vee\ r \Leftrightarrow \neg\mathsf{distincts}(\ldots)$, which
\texttt{-split} does not split, and in the context of the second disjunct neither literal of the negated
first is available, so the condition is not decided; probes (rules with the real left-hand side and the
first $k$ conditions) located it as the only failing one. As two implications, each goal has its literal
of the exit, $r$ respectively $n \le i$ with $i = \intf(n')$, in the context of the atom.

An earlier set of lemmas closed both goals and also the negated postcondition: the rule for the first step
had the summary fragment as an open side condition, and a rule $\mathsf{nodup}(a, k, n) \Leftrightarrow \mathsf{true}$
under the hypothesis $\mathsf{distincts}(\ldots)$ for the second. An $\mathsf{all}$ side condition is matched
in the simplifier's context of the fact's combinator, where rules fire: under the hypothesis $\mathsf{distincts}$,
the rule rewrote the condition to \eqref{eq:fd1}, a fact, so $\mathsf{distincts}$ was derived from itself
(finding 13 of \texttt{NOTES.md}). The rules above take the atom $\mathsf{distincts}$ as left-hand side
and never rewrite $\mathsf{nodup}$.

\subsection{The idea}
An early exit is a data-dependent number of iterations with a disjunctive exit; the summary of the
inner loop must be stated so that it holds in both cases, which the abstract atom $\mathsf{nodup}$
affords, since it describes the result rather than the number of iterations; the outer loop is summarized
on the atom of the postcondition with one rule per case of its exit, each with the case as a closed
condition. The soundness lesson stands: a side condition that other rules can simplify under the conclusion
to be proved is a circular justification, and the safe form of side conditions is an atom that is matched
literally (as \lemma{consts\_prefix\_def} of the database does for a constraint, and as the summary itself
does for the inner loop).

\section{Binary Search: the Lower Bound (open)}\label{sec:lb}

\subsection{The program}
\begin{lstlisting}
method main(a: array<int, int>, n: int, x: int) returns (lo: int)
  requires 0 <= n
  requires sorted(a, 0, n)
  ensures 0 <= lo <= n
  ensures forall k: int :: 0 <= k < lo ==> a[k] < x
  ensures forall k: int :: lo <= k < n ==> x <= a[k]
{
  var hi: int; var mid: int;
  lo := 0; hi := n;
  while lo < hi {
    mid := (lo + hi) / 2;
    if a[mid] < x { lo := mid + 1; } else { hi := mid; }
  }
}
\end{lstlisting}
The first position of $a[0..n)$ whose element is at least $x$, for a sorted array. The example
\texttt{bsearch.dfy} of \texttt{examples/dafny/teq} verifies a binary search with the stronger
precondition of partition points for all indices; here the precondition is $\sorted(a, 0, n)$, so the
comparisons $a[m] < x$ are meaningful for $0 \le m < n$ only.

\subsection{The trace equation system}
\begin{align}
  K &= \cnt(0, 1, n') \qquad lo = \cur(L) \qquad hi = \cur(H) \qquad hi \le lo \label{eq:lb0}\\
  L &= \tr{0}{\bl{K}, \bl{H}}{lo, k, hi \mapsto \upd{G}{(hi + lo)/2 + 1}\ \ \upd{\neg G}{lo}} \label{eq:lb1}\\
  H &= \tr{n}{\bl{K}, \bl{L}}{hi, k, lo \mapsto \upd{G}{hi}\ \ \upd{\neg G}{(hi + lo)/2}} \label{eq:lb2}\\
  &\qquad\text{where } G \text{ abbreviates } a[(hi + lo)/2] < x \nonumber\\
  &\al{\bl{K}, \bl{L}, \bl{H}}{k, lo, hi \mapsto lo < hi} \label{eq:lb3}
\end{align}
As in the merge, the two variables are traces over each other, with a data-dependent guard that reads the
array at a position computed from both.

\subsection{The attempt}
The intended lemmas name the result by an abstract function $\mathsf{lb}(a, n, x)$ with the facts
$0 \le \mathsf{lb}(a, n, x) \le n$, $(a[k] < x) \Leftrightarrow \mathsf{true}$ for $0 \le k < \mathsf{lb}(a, n, x)$,
and $(x \le a[k]) \Leftrightarrow \mathsf{true}$ for $\mathsf{lb}(a, n, x) \le k < n$ (all under $\sorted(a, 0, n)$),
which close the postcondition once $lo = \mathsf{lb}(a, n, x)$ is known. The summary of the loop in the style of
the merge (Section~\ref{sec:merge}): the guard is the higher-order $G(lo, hi)$ pinned by
$\forall lo, hi.\ G(lo, hi) \Leftrightarrow a[(hi + lo)/2] < x$, the trace of $hi$ is tied by an open trace variable
$Hv = \tr{n}{\bl{\cnt(0, 1, nn)}, \bl{Lv}}{\ldots}$ with $hs = \bl{Hv}$, and the left-hand side is
$\cur(\tr{0}{\bl{\cnt(0, 1, nn)}, hs}{lo, k, hi \mapsto \upd{G(lo, hi)}{(hi + lo)/2 + 1}\ \ \upd{\neg G(lo, hi)}{lo}})$,
rewritten to $\mathsf{lb}(a, n, x)$. What the lemma must also know is that the loop has terminated, the exit
$hi \le lo$ of \eqref{eq:lb0}: without it the current value of $L$ is not the lower bound.

Stating the exit as a side condition $\cur(Hv) \le \cur(Lv)$ makes the run exceed the time limit: to decide
the condition, the simplifier rewrites $\cur(Lv)$, which is the left-hand side of the very lemma, whose
conditions decide it again (finding 20 of \texttt{NOTES.md}). Stating it as a comparison $h \le l$ of two
open variables, matched against the hypothesis $hi \le lo$, makes the lemma fire, and the four facts about
$\mathsf{lb}$ then close the three postconditions: the goal is closed, and a wrong postcondition leaves a
goal. But $h$ and $l$ are not tied to the traces, so the rule is not a sound lemma (in this goal its
conclusion holds whichever comparison the condition matched, since the exit does hold in the final state).
Every sound tie fails for a reason of the engine, located with probes: the instance of a side condition is
simplified before it is decided or matched, so $l = \cur(Lv)$ with $Lv$ bound re-enters the lemma like the
first form, whether closed or open. With $Lv$ open, $l = \cur(Lv)$ binds both from the hypothesis $lo = \cur(L)$
without re-entrance; but then the left-hand side, matched against $\cur(L)$ through the definition of the
trace variable $L$, turns out to have left its dependency variables $nn$ and $hs$ unbound (a probe fired with
"variables of a rule not instantiated"), and the trace condition that would bind them,
$Lv = \tr{0}{\bl{\cnt(0, 1, nn)}, hs}{\ldots}$, no longer matches once $Lv$ is bound, since its instance is
simplified with $L$ replaced by its definition.

\subsection{The idea, and what is missing}
Binary search is the paper's acknowledged weakness: the loop computes a bound by narrowing, and only a
partial description of the traces (monotone $lo$, monotone $hi$, $lo \le \mathsf{lb} \le hi$ throughout) is
needed, not their values. \texttt{bsearch.dfy} handles it by abstracting the comparisons with the array
into comparisons of indices with partition points (lemmas $(a[m] < x) \Leftrightarrow m < p$ without side
conditions), by the named monotone traces $\mathsf{incs}$ and $\mathsf{decs}$ whose bounds are their
facts, and by a cut after the loop. With $\sorted$ instead of partition points, the comparison lemma needs
$0 \le m < n$, which in the context of the updates follows from the invariant $0 \le lo \le hi \le n$,
the thing to be established; the summary lemma of this section bypasses the invariant by stating the
result directly, and founders on tying the exit to the traces soundly. What is missing is in the engine: a
rule should not be applied while the instances of its own side conditions are simplified, and a left-hand
side matched through a trace variable should bind all its variables; with either, the sound form of the
lemma closes the goal.

\section{Array Programs in C: the MonoCera Tests}\label{sec:monocera}

The MonoCera evaluation (Amilon, Esen, Gurov, Lidström, Rümmer, \emph{Automatic Program Instrumentation for
Automatic Verification}, CAV 2023) has 125 C programs derived from SV-COMP, each with a loop that fills an array and
a loop that checks an assertion over it, one test per aggregation that the assertion admits ($\forall$, $\max$,
$\min$, $\sum$). Unlike the Dafny programs of the previous sections, they carry no specification and no lemmas: the
assertion is the postcondition, and everything a proof needs must come from the database. They are run with all
modules of the database and mathematical integers, as in that evaluation:
\texttt{-unbounded -teq} with the modules, \texttt{-prove -split -prove -split -prove}. The second and third rounds
split the disjunctions that the first introduces, e.g.\ the case of no iteration.

\begin{table}[t]
\centering\small
\begin{tabular}{@{}lrrrrr@{}}
\toprule
aggregation & tests & lookup (\texttt{-prove}) & lookup and Z3 & MonoCera & Horn clauses for Z3 \\
\midrule
$\forall$ & 93 & 53 & 59 & 26 & 0 \\
$\max$ & 4 & 2 & 3 & 2 & 0 \\
$\min$ & 11 & 7 & 9 & 3 & 0 \\
$\sum$ & 17 & 17 & 17 & 10 & 0 \\
\midrule
total & 125 & 79 & 88 & 41 & 0 \\
\bottomrule
\end{tabular}
\caption{Proved tests of the MonoCera evaluation, with mathematical integers: lookup by the database alone, lookup with
the remaining goals lowered for Z3 (\texttt{-fallback}), MonoCera, the instrumenting verifier of that evaluation, and
Horn clauses for Z3 from the same frontend. No result is wrong. The lemmas for these tests started in a separate file
(51 tests in a first run, 73 with that file); their generic part has since moved into the database, and the
specific lemmas (strides, stores ahead) are in \texttt{benchmarks/monocera/lemmas.dfy}, which these runs do not
use.}\label{tab:monocera}
\end{table}


Two things distinguish the C programs. A \texttt{for} loop with a body of conditionals gives one case per path, also
where the paths agree, e.g.\ the counter $i$ with two cases $i + 1$, which the database merges first
(\lemma{trace\_merge\_cases\_ne}). And the harness branches, \texttt{if (N <= 0) return 1}, so the goal has one
branch per path, of which one carries the loops. The iteration constraint of a loop with guarded cases also states
that one of the guards held, e.g.\ $i \bmod 2 = 0 \vee \neg(i \bmod 2 = 0)$: the completeness of the guards in the
context of the update (Section~\ref{sec:engine:cases}) is a fact of the goal.

\subsection{Loops in sequence: sums of filled arrays}\label{sec:brs2}
\begin{lstlisting}[language=C]
for (i = 0; i < N; i++) { if (i % 2 == 0) a[i] = 2; else a[i] = 0; }   // brs2
for (i = 0; i < N; i++) { if (i == 0) sum[0] = 0; else sum[0] = sum[0] + a[i]; }
__VERIFIER_assert(sum[0] <= 2 * N);
\end{lstlisting}
The second loop reads the final array of the first, and the sum is kept in the cell $0$ of a one-element array, so
its trace is a trace of stores at the constant key $0$. The proof needs one lemma per shape, each landing in the
normal form that the next one expects:
\begin{enumerate}[nosep]
\item \emph{The counters.} The traces of $i$ are counters, the constraints and exits give $\intf(n) = N$ for both
  loops, and inside each iteration $i$ is the iteration $k$.
\item \emph{The first loop is stores of a map.} Its update stores a constant chosen by a condition on the index:
\begin{align*}
  &\lemma{stores\_index\_cases\_def}:\\
  &\quad \mathsf{trace}(a_0,\ \mathsf{for}\ \bl{\cnt(0, 1, n)}, \bl{\cnt(0, 1, n)}\\
  &\qquad\quad \mathsf{do}\ a, k, i \mapsto \upd{P(k)}{a[k := c]}\ \ \upd{\neg P(k)}{a[k := d]})\\
  &\quad = \mathsf{stores}(a_0, \bl{\cnt(0, 1, n)}, \mapc{\bl{\cnt(0, 1, n)}}{k \mapsto \ite{P(k)}{c}{d}})
\end{align*}
  with $P$ the parity, abstracted from the guard.
\item \emph{The read in the second loop.} $a'[k]$ is the stored value at $k$ (\lemma{current\_stores}, with the
  membership and position of $k$ in the counter, \lemma{occurs\_counter\_up} and \lemma{indexof\_counter\_up}), decided
  in the context of the iteration, where $0 \le k < \intf(n)$.
\item \emph{The cell is a sum.} The trace of stores into the cell is read as the scalar sum:
\begin{align*}
  &\lemma{select\_current\_trace\_reset\_add}:\\
  &\quad \cur(\mathsf{trace}(a_0, \ldots\ \upd{k = 0}{a[j := c]}\ \ \upd{0 < k}{a[j := a[j] + G(k)]}))[j]\\
  &\quad = \ite{n = 0}{a_0[j]}{c + \mathsf{sum}(\mapc{\bl{\cnt(0, 1, n)}}{k \mapsto \ite{0 < k}{G(k)}{0}})}
\end{align*}
\item \emph{The bound.} \lemma{sum\_cases\_below} bounds a sum of terms that are each $0$, $c$ or $d$ by $e$; its
  left-hand side is the comparison $\mathsf{sum}(\ldots) \le e$, whose right side binds $e$, so the bound need not be
  a fact.
\end{enumerate}
The test \texttt{sum/ms2}, which stores $i \bmod 2$, differs only in the bound (\lemma{sum\_mod\_below}), and
\texttt{sum/conda} only in the first loops: an array filled with ones and then updated in place by cases on the
old value,
\begin{align*}
  &\lemma{stores\_cases\_def}:\\
  &\quad \tr{a_0}{\ldots}{a, k, i \mapsto \upd{P(a[k])}{a[k := c]}\ \ \upd{\neg P(a[k])}{a[k := G(a[k])]}}\\
  &\quad = \mathsf{stores}(a_0, \ldots, \mapc{\ldots}{k \mapsto \ite{P(a_0[k])}{c}{G(a_0[k])}}),
\end{align*}
since each position is read only in its own iteration, before it is written. Its first case stores a constant, not
$G'(a[k])$: by the time the lemma is tried, the read $a[k]$ inside the update has already become $1$ (the stores of
the previous loop, read in the context of the iteration), so the program's $a[k] + 1$ is $2$. A higher-order
instance must contain the term it abstracts, and the constant $2$ contains no $a[k]$.

\paragraph{What this needs.} Loops in sequence compose through their current values: each loop reads the previous
array, the read is resolved in the context of the iteration, and only then is the shape of the next loop recognized.
The lemmas are written for the goal as the engine has simplified it, not for the program text. And where the data has
a different shape, a cell of an array instead of a scalar, the lemma for the new shape lands in the normal form of the
old one (the right-hand side of the scalar sum), so that everything after it is shared. One detail of the engine is
needed here: the instance $(z) \mapsto a[z]$ of a higher-order variable must not be eta-reduced to the array $a$,
since $a(k)$ applied as a function is not the read $a[k]$ that the rules are about; eta reduction applies only to
functions of array type.

\subsection{Strided writes: a quotient in the bound (open)}\label{sec:pr2}
\begin{lstlisting}[language=C]
if (CELLCOUNT % 2 != 0) return 1;   // pr2
for (i = 1; i <= CELLCOUNT / 2; i++) {
  volArray[i*2 - 2] = (3 >= MINVAL) ? 3 : 0;
  volArray[i*2 - 1] = (1 >= MINVAL) ? 1 : 0;
}
for (i = 0; i < CELLCOUNT; i++)
  __VERIFIER_assert(volArray[i] >= MINVAL || volArray[i] == 0);
\end{lstlisting}
The loop stores pairs at $2k$ and $2k + 1$, with values chosen by three cases on \texttt{MINVAL} that do not depend
on the iteration. The module for strided stores reads a cell of the result:
\begin{align*}
  &\lemma{select\_current\_trace\_stride2\_cases}:\quad R \Leftrightarrow (\neg P \wedge \neg Q) \Longrightarrow\\
  &\quad \cur(\mathsf{trace}(a_0, \ldots\ \upd{P}{a[2k := X_1][2k + 1 := Y_1]}\ \ \upd{Q}{\ldots}\ \ \upd{R}{\ldots}))[j]\\
  &\quad = \ite{0 \le j < 2\,\intf(n)}{\ldots}{a_0[j]},
\end{align*}
by the parity of $j$ and the guards. Its side condition requires the guards to be exhaustive; the lemma fixes two
stores and three cases on its left-hand side. It fires on the reads of the checking loop, and every value it
produces satisfies the assertion. The goal that remains is the range above the stores:
\[ 2\,(C / 2) < C \ \wedge\ C \bmod 2 = 0 \ \wedge\ 2\,(C / 2) = i \Longrightarrow \ldots a[i] \ldots \]
for $C$ the size: the loop ran $C / 2$ times, the read at $i = 2\,(C / 2)$ lies beyond the stores, and the case is
impossible since $C$ is even. The arithmetic knows the bounds of $C / 2$, but not that $C \bmod 2$ is
$C - 2\,(C / 2)$, so it does not see the contradiction.

\paragraph{What this needs.} Remainders as linear terms over quotients: with $t \bmod k = t - k \cdot (t / k)$ for a
positive constant $k$, the facts $C \bmod 2 = 0$ and $2\,(C/2) < C$ are inconsistent, and the test closes. A lemma
stating the fact for one form of the equation, e.g.\ $x \bmod 2 = 0 \Rightarrow (n = x / 2 \Leftrightarrow x = 2n)$,
depends on the form in which the engine keeps the equation, which changes with every improvement of the arithmetic;
the remainder belongs to the decision procedure. The tests with $s$ per-cell guards, \texttt{pr3} to \texttt{pr5},
need $2^s$ combinations of cases, which a lemma with a fixed number of cases does not provide: a lemma schema with
the number of stores and cases as a parameter.

\subsection{Stores ahead, in both directions}\label{sec:rewnif}
\begin{lstlisting}[language=C]
for (i = 0; i < SIZE; i++) { if (i + 1 < SIZE) a[i+1] = i; a[i] = i; }       // rewnif
for (i = SIZE - 1; i >= 0; i--) { if (i - 1 >= 0) a[i-1] = i - 2; a[i] = i; } // rewnifrev
for (i = 0; i < SIZE; i++) __VERIFIER_assert(a[i] >= i);
\end{lstlisting}
Each iteration writes one cell ahead before it writes its own; the value written ahead is overwritten by the next
iteration, so a cell holds what its own iteration wrote last, except the cell past the last iteration, which holds
the value written ahead. The module for strided stores has one lemma per direction; for the loop down from $s$,
\begin{lstlisting}
lemma {:rw} {:teq} select_current_trace_ahead_down(a0: array<int, int>, n: Nat, s: int,
                                                   F: (int) -> int, G: (int) -> int, j: int)
  ensures
    current(trace(a0, for ⌊counter(0, 1, n)⌋, ⌊counter(s, -1, n)⌋ do a, k, i ::
              case 1 <= i => a[i + (-1) := F(i)][i := G(i)]
              case i < 1 => a[i := G(i)]))[j]
      == if s - integer(n) < j && j <= s then G(j)
         else if j == s - integer(n) && 0 <= j && 0 < integer(n) then F(j + 1)
         else a0[j]
\end{lstlisting}
It fires on the reads of the checking loop with $s := \texttt{SIZE} - 1$, $F := (z) \mapsto z - 2$ and
$G := (z) \mapsto z$, and the assertion holds in every branch; the test closes under $\forall$ and $\min$.

\paragraph{What this needs.} The pattern contains two literals that a matcher must treat modulo the normal form of the
goal: the step $-1$ of the counter dependency, matched through the definition of the program's counter, and the store
position $i + (-1)$, which the goal has as $i - 1$, a sum with a negative literal. The engine compares a negative
literal among the summands as the negation of its absolute value, so the lemma can be written as the program reads.
Without that, each such position needs a higher-order variable $P(i)$ pinned to $i - 1$ by a side condition, and each
literal step a variable $d$ with $d = -1$. The third variant of the test, \texttt{rewnifrev2}, writes ahead upwards in a
loop that runs downwards, a shape for which there is no lemma.

\subsection{Recurrences on the previous cell}\label{sec:ifeqn2}
\begin{lstlisting}[language=C]
for (i = 0; i < N; i++) { if (i == 0) a[0] = 4; else a[i] = a[i-1] + 4; }   // ifeqn2
for (i = 0; i < N; i++) { if (i == 0) b[0] = 1; else b[i] = b[i-1] + a[i-1]; }
for (i = 0; i < N; i++) __VERIFIER_assert(b[i] == 2*i*i + 2*i + 1);
\end{lstlisting}
A store at $k$ of a value read at $k - 1$ is not a store of a map: the value depends on the trace, not on the initial
array. The aggregation module names the recurrence by its closed form,
\begin{lstlisting}
lemma {:rw} {:teq} select_current_trace_recurrence(a0: array<int, int>, n: Nat, c: int, j: int,
                                                   G: (int) -> int)
  ensures current(trace(a0, for ⌊counter(0, 1, n)⌋ do a, i ::
            case i == 0 => a[0 := c] case 0 < i => a[i := a[i - 1] + G(i)]))[j]
    == if 0 <= j && j < integer(n)
       then c + sum(maps(for ⌊counter(1, 1, natural(j))⌋ yield k :: G(k)))
       else a0[j]
\end{lstlisting}
The first loop gives $a'[j] = 4 + \mathsf{sum}(\mapc{\bl{\cnt(1, 1, \natf(j))}}{k \mapsto 4}) = 4 + 4j$ (the map of a
constant and its sum). In the second loop the summand $G(k) = a'[k - 1]$ is then $4k$, so
$b'[j] = 1 + \mathsf{sum}(\mapc{\bl{\cnt(1, 1, \natf(j))}}{k \mapsto 4k})$; the map of a multiple of the counter is
a counter, $\bl{\cnt(4, 4, \natf(j))}$ (\lemma{maps\_counter\_scale}), whose sum is the arithmetic series
(\lemma{sum\_counter}), $2j^2 + 2j$. The read of the checking loop at its index meets the assertion exactly.

\paragraph{What this needs.} The read $a[i - 1]$ of the pattern must match the goal's $a[k + (-1)]$, the same negative
literal as above; the higher-order summand $G$ must take the read of the other array as a whole; and the closed forms
of maps and sums over counters must reach polynomials, which the arithmetic compares. The recurrences with other
increments, products or conditions (\texttt{ifeqn3} to \texttt{ifeqn5}, \texttt{sqm}, \texttt{condm}) need their own
closed forms; a recurrence solver as a generator of such lemmas would cover the polynomial ones.

\section{Alternating Values: a Sum Bounded Through Three Lemmas}\label{sec:brs2}

\subsection{The program}
The test \texttt{sum/brs2}: the array alternates between $2$ and $0$, the sum is kept in the cell $0$ of a
one-element array, and the assertion is an upper bound.
\begin{lstlisting}[language=C]
for (i = 0; i < N; i++) { if (i % 2 == 0) a[i] = 2; else a[i] = 0; }
for (i = 0; i < N; i++) { if (i == 0) sum[0] = 0; else sum[0] = sum[0] + a[i]; }
__VERIFIER_assert(sum[0] <= 2 * N);
\end{lstlisting}

\subsection{The trace equation system}
Two loops in sequence; the second one's trace of \texttt{sum} has the final array $a'$ of the first in its
update. With $N > 0$ from the harness,
\begin{align}
  K &= \cnt(0, 1, n) \qquad a' = \cur(A) \qquad i' = \cur(I) \label{eq:brs0}\\
  A &= \tr{a_0}{\bl{K}, \bl{I}}{a, k, i \mapsto \upd{i \bmod 2 = 0}{a[i := 2]}\ \ \upd{i \bmod 2 \ne 0}{a[i := 0]}} \label{eq:brs1}\\
  I &= \tr{0}{\bl{K}}{i, k \mapsto \upd{i \bmod 2 = 0}{i + 1}\ \ \upd{i \bmod 2 \ne 0}{i + 1}} \label{eq:brs2}\\
  &\al{\bl{K}, \bl{A}, \bl{I}}{k, a, i \mapsto i < N} \qquad \neg(i' < N) \label{eq:brs3}\\
  K' &= \cnt(0, 1, n') \qquad s' = \cur(S) \qquad i'' = \cur(I') \nonumber\\
  S &= \tr{s_0}{\bl{K'}, \bl{I'}}{s, k, i \mapsto \upd{i = 0}{s[0 := 0]}\ \ \upd{i \ne 0}{s[0 := s[0] + a'[i]]}} \label{eq:brs4}\\
  I' &= \tr{0}{\bl{K'}}{i, k \mapsto i + 1} \qquad \al{\bl{K'}, \bl{S}, \bl{I'}}{k, s, i \mapsto i < N} \qquad \neg(i'' < N) \nonumber\\
  &\Longrightarrow\ s'[0] \le 2 N \nonumber
\end{align}
(the \texttt{rem} of C is the remainder of the prelude, which the simplifier turns into $\bmod$ under
$0 \le i$). The accumulator is not a scalar but a cell of an array, so \eqref{eq:brs4} is a trace of stores
at the constant key $0$, and the array it reads is the current value of another trace of stores.

\subsection{The lemmas}
Three lemmas of the database cooperate, in the order in which the goal needs them.
\begin{align*}
  &\lemma{stores\_index\_cases\_def}:\\
  &\quad \tr{a_0}{\bl{\cnt(0, 1, n)}, \bl{\cnt(0, 1, n)}}{a, k, i \mapsto \upd{P(k)}{a[k := c]}\ \ \upd{\neg P(k)}{a[k := d]}}\\
  &\qquad = \mathsf{stores}(a_0, \bl{\cnt(0, 1, n)}, \mapc{\bl{\cnt(0, 1, n)}}{k \mapsto \ite{P(k)}{c}{d}})\\[4pt]
  &\lemma{select\_current\_trace\_reset\_add}:\\
  &\quad \cur(\tr{a_0}{\bl{\cnt(0, 1, n)}, \bl{\cnt(0, 1, n)}}{a, k, i \mapsto \upd{k = 0}{a[j := c]}\ \ \upd{0 < k}{a[j := a[j] + G(k)]}})[j]\\
  &\qquad = \ite{n = 0}{a_0[j]}{c + \mathsf{sum}(\mapc{\bl{\cnt(0, 1, n)}}{k \mapsto \ite{0 < k}{G(k)}{0}})}\\[4pt]
  &\lemma{sum\_cases\_below}:\quad 0 \le e\ \wedge\ \intf(n) \cdot c \le e\ \wedge\ \intf(n) \cdot d \le e\\
  &\quad\Longrightarrow\ \mathsf{sum}(\mapc{\bl{\cnt(0, 1, n)}}{k \mapsto \ite{0 < k}{(\ite{P(k)}{c}{d})}{0}}) \le e\ \Longleftrightarrow\ \mathsf{true}
\end{align*}
The first turns a loop that stores at each position a constant chosen by cases on the \emph{index} into
stores of a map; the second reads a sum kept in one cell $j$ of an array as the scalar sum, with the
right-hand side of \lemma{current\_trace\_reset\_add} for a scalar accumulator, so that the database continues
from there; the third bounds a sum of $n$ terms that are each $0$, $c$ or $d$, with its left-hand side a
comparison whose constants bind the bound $e$ (the general \lemma{sum\_below\_prefix} of the database does not
apply, since its bound occurs only in a side condition, which would have to be a fact).

\subsection{The proof}
Of the 62 rule applications, those that carry the proof:
\begin{enumerate}[nosep]
\item \textbf{The counters.} \lemma{trace\_merge\_cases\_ne} joins the two cases of \eqref{eq:brs2}, both
  $i + 1$, and \lemma{counter\_def} makes $I = \cnt(0, 1, n)$, likewise $I' = \cnt(0, 1, n')$. In the
  context of the iteration $k$ the simplifier reads $i$ as $k$ (\lemma{earlier\_counter}), so the constraint
  \eqref{eq:brs3} is $\al{\bl{\cnt(0, 1, n)}}{k \mapsto k < N}$, which \lemma{invs\_counter\_upper\_nonnegative}
  resolves to $\intf(n) \le N$; with the exit $N \le i' = \intf(n)$ (\lemma{current\_counter}) this gives
  $\intf(n) = N$, and for the second loop $\intf(n') = N$, hence $n' = n$ (\lemma{integer\_injective}).
\item \textbf{The first loop is stores.} In the update of \eqref{eq:brs1}, $i \bmod 2 = 0$ is $k \bmod 2 = 0$
  and the guards are $k \bmod 2 = 0$ and $0 < k \bmod 2$; \lemma{stores\_index\_cases\_def} matches with
  $P := \lambda k.\ k \bmod 2 = 0$, $c := 2$, $d := 0$ (the higher-order $P$ by abstraction of the guard), so
  $a' = \cur(\mathsf{stores}(a_0, \bl{\cnt(0, 1, n)}, \mapc{\bl{\cnt(0, 1, n)}}{k \mapsto \ite{k \bmod 2 = 0}{2}{0}}))$.
\item \textbf{The read in the second loop.} The update of \eqref{eq:brs4} reads $a'[i]$, i.e.\ $a'[k]$;
  \lemma{current\_stores} makes it $\ite{\mathsf{occurs}(k, \bl{\cnt(0, 1, n)})}{\mathsf{earlier}(\mathsf{indexof}(k, \bl{\cnt(0, 1, n)}), \mapc{\ldots}{\ldots})}{a_0[k]}$,
  \lemma{occurs\_counter\_up} the test $0 \le k < \intf(n)$, which the context of the iteration decides (it
  has $0 \le k < n' = n$), \lemma{indexof\_counter\_up} the index $\natf(k)$, and the read of the map at $k$ is
  $\ite{k \bmod 2 = 0}{2}{0}$. The second loop is now
  $\tr{s_0}{\ldots}{s, k, i \mapsto \upd{k = 0}{s[0 := 0]}\ \ \upd{0 < k}{s[0 := s[0] + (\ite{k \bmod 2 = 0}{2}{0})]}}$:
  a sum in a cell, with the summand a function of the iteration.
\item \textbf{The cell is a sum.} The postcondition reads $s'[0] = \cur(S)[0]$;
  \lemma{select\_current\_trace\_reset\_add} matches with $j := 0$, $c := 0$ and
  $G := \lambda k.\ \ite{k \bmod 2 = 0}{2}{0}$, abstracted from the summand, and yields
  $\ite{n = 0}{s_0[0]}{0 + \mathsf{sum}(\mapc{\bl{\cnt(0, 1, n)}}{k \mapsto \ite{0 < k}{G(k)}{0}})}$; the
  case $n = 0$ is $\intf(n) = 0$ (\lemma{zero\_integer}), false under $\intf(n) = N > 0$.
\item \textbf{The bound.} \lemma{sum\_cases\_below} matches the remaining atom
  $\mathsf{sum}(\ldots) \le 2N$ with $e := 2N$, $c := 2$, $d := 0$ and $P$ the parity; its conditions
  $0 \le 2N$, $2\,\intf(n) \le 2N$ and $0 \le 2N$ hold by $\intf(n) = N$. The goal is closed.
\end{enumerate}
A variant asserting \texttt{sum[0] <= N - 2} stays open with the residual $\neg(0 < N)$: the lemmas compute
the value $N - 1$ of the sum (the summand $2$ at the even positions from $2$ on) and the bound is false; the
same lemmas close \texttt{brs3} to \texttt{brs5}, and \lemma{sum\_mod\_below} in place of the third closes
\texttt{ms2} to \texttt{ms5}, whose arrays hold $i \bmod m$ (Section~\ref{sec:ms2}).

\subsection{The idea}
The scalar version of this program, which the Dafny translations of SV-COMP produce, is closed by
\lemma{current\_trace\_reset\_add} alone; the C program needs the same lemma for a cell of an array, with
the same right-hand side. This is ``name the normal form'' of Section~\ref{sec:principles:shape} at the level
of the database: the shape of the data decides which lemma matches, and the lemma for the new shape
should land in the normal form of the old, so that everything after it is shared. The bug that this test
found on the way is finding~10 of the algorithms seen from the engine: the simplifier's eta reduction had
turned the instance $\lambda z.\ a[z]$ of $G$ into the array $a$ itself, which left an application $a(k)$
that no rule about $a[k]$ matched, fixed in version 0.35.15 by restricting eta reduction to lambdas of
array type.

\section{An Update in Place and a Sum of Remainders}\label{sec:conda}

\subsection{The program}
The test \texttt{sum/conda} fills an array with ones, modifies it in place by cases on the old value, and
sums it; the assertion is the exact sum.
\begin{lstlisting}[language=C]
sum[0] = 0;
for (i = 0; i < N; i++) a[i] = 1;
for (i = 0; i < N; i++) { if (a[i] == 1) a[i] = a[i] + 1; else a[i] = a[i] - 1; }
for (i = 0; i < N; i++) sum[0] = sum[0] + a[i];
__VERIFIER_assert(sum[0] == 2 * N);
\end{lstlisting}

\subsection{The trace equation system}
Three loops in sequence, each over its own counter of $n$, $n'$, $n''$ steps, all equal to $N$ by their
tests and exits as in Section~\ref{sec:brs2}. The middle one is the interesting equation: its update reads
the cell it writes, and its counter trace $I'$ has the array among its dependencies, since the guards
mention it,
\begin{align}
  A &= \tr{a_0}{\bl{K}, \bl{I}}{a, k, i \mapsto a[i := 1]} \label{eq:cd1}\\
  A' &= \tr{a'}{\bl{K'}, \bl{I'}}{a, k, i \mapsto \upd{a[i] = 1}{a[i := a[i] + 1]}\ \ \upd{a[i] \ne 1}{a[i := a[i] - 1]}} \label{eq:cd2}\\
  I' &= \tr{0}{\bl{K'}, \bl{A'}}{i, k, a \mapsto \upd{a[i] = 1}{i + 1}\ \ \upd{a[i] \ne 1}{i + 1}} \nonumber\\
  S &= \tr{s_0[0 := 0]}{\bl{K''}, \bl{I''}}{s, k, i \mapsto s[0 := s[0] + a''[i]]} \label{eq:cd3}\\
  &\Longrightarrow\ \cur(S)[0] = 2N \nonumber
\end{align}
with $a' = \cur(A)$ and $a'' = \cur(A')$.

\subsection{The lemmas}
\begin{align*}
  &\lemma{stores\_cases\_def}:\\
  &\quad \mathsf{trace}(a_0, \mathsf{for}\ \bl{\cnt(0, 1, n)}, \bl{\cnt(0, 1, n)}\ \mathsf{do}\ a, k, i \mapsto {}\\
  &\qquad\qquad \upd{P(a[k])}{a[k := c]}\ \ \upd{\neg P(a[k])}{a[k := G(a[k])]})\\
  &\quad = \mathsf{stores}(a_0, \bl{\cnt(0, 1, n)}, \mapc{\bl{\cnt(0, 1, n)}}{k \mapsto \ite{P(a_0[k])}{c}{G(a_0[k])}})
\end{align*}
Each position $k$ is written once, in iteration $k$, and read only there before it is written, so the loop
stores values computed from the \emph{initial} array. The first case stores a constant and not $G'(a[k])$: a
higher-order instance must contain the term it abstracts (\texttt{Rewrite.abstraction}), and $a[k] + 1$ is
$2$ in this program by the time the lemma is tried (step~2 below). The rest is the database's:
\lemma{stores\_def} and \lemma{maps\_const} for the first loop, \lemma{current\_stores} with
\lemma{occurs\_counter\_up}, \lemma{indexof\_counter\_up} and \lemma{earlier\_consts} for the reads,
\lemma{current\_trace\_add\_at} for the sum at the constant key, \lemma{selects\_stores\_counter} and
\lemma{sum\_consts} for the result.

\subsection{The proof}
Of the 69 rule applications:
\begin{enumerate}[nosep]
\item \textbf{The first loop.} \eqref{eq:cd1} is $\mathsf{stores}(a_0, \bl{I}, \mapc{\bl{K}}{k \mapsto 1})$ by
  \lemma{stores\_def}, and the map of a constant is $\bl{\mathsf{consts}(n, 1)}$ (\lemma{maps\_const}); $I$
  is $\cnt(0, 1, n)$.
\item \textbf{The guards of the second loop.} In the context of the iteration $k$ of \eqref{eq:cd2}, the
  previous value $a$ at position $k$ is unmodified, $a[k] = a'[k]$, and $a'[k]$ is read from the stores of the
  first loop: \lemma{current\_stores} gives
  $\ite{\mathsf{occurs}(k, \bl{I})}{\mathsf{earlier}(\mathsf{indexof}(k, \bl{I}), \bl{\mathsf{consts}(n, 1)})}{a_0[k]}$,
  the test is $0 \le k < \intf(n)$ by \lemma{occurs\_counter\_up}, true in the iteration, and the read of the
  constant sequence is $1$ (\lemma{earlier\_consts}). So $a[k] = 1$ inside the update: the first result
  $a[k] + 1$ is $2$, the guards are $a[k] = 1$ and $a[k] \ne 1$ on the \emph{bound} variable $a$, which the
  lemma's $P$ abstracts.
\item \textbf{The second loop is stores.} \lemma{stores\_cases\_def} matches \eqref{eq:cd2} with
  $P := \lambda z.\ z = 1$, $c := 2$, $G := \lambda z.\ z - 1$ and gives
  $\mathsf{stores}(a', \bl{\cnt(0, 1, n')}, \mapc{\bl{\cnt(0, 1, n')}}{k \mapsto \ite{a'[k] = 1}{2}{a'[k] - 1}})$;
  the read $a'[k]$ in the map is again $1$ by step~2, the conditional is $2$, and the map is
  $\bl{\mathsf{consts}(n', 2)}$.
\item \textbf{The sum.} \eqref{eq:cd3} is a trace that adds to the cell $0$ in every iteration;
  \lemma{current\_trace\_add\_at} makes $\cur(S) = s_0[0 := 0][0 := 0 + \mathsf{sum}(\mapc{\bl{K''}, \bl{I''}}{k, i \mapsto a''[i]})]$,
  the map of reads is $\mathsf{selects}(a'', \bl{I''})$ (\lemma{selects}), and with $n'' = n' = n$
  (\lemma{integer\_injective} from the three exits) \lemma{selects\_stores\_counter} reads the stored values
  back at the positions of the counter: $\bl{\mathsf{consts}(n, 2)}$, whose sum is $2\,\intf(n)$
  (\lemma{sum\_consts}). The postcondition $2\,\intf(n) = 2N$ holds by $\intf(n) = N$.
\end{enumerate}
The tests \texttt{s1lif} to \texttt{s5lif}, whose second loops have the same shape with other cases, close by
the same steps (69 rule applications for \texttt{s1lif}); a false assertion \texttt{sum[0] == 2 * N + 1}
leaves $2\,\intf(n) = 2N + 1$, i.e.\ the lemmas compute the sum rather than close the goal.

\subsection{A sum of remainders}\label{sec:ms2}
The test \texttt{sum/ms2} stores $i \bmod 2$ at $i$ and asserts \texttt{sum[0] <= N} for the sum of
Section~\ref{sec:brs2}. The first loop is plain stores of a map, $\mathsf{stores}(a_0, \bl{I}, \mapc{\bl{K}, \bl{I}}{k, i \mapsto \ite{0 \le i}{i \bmod 2}{-(-i \bmod 2)}})$
(the remainder of C under a sign that the context of the map does not know yet), the read $a'[k]$ in the
second loop goes through \lemma{current\_stores} to the map at $k$, where $0 \le k$ holds and the read is
$k \bmod 2$, \lemma{select\_current\_trace\_reset\_add} abstracts $G := \lambda k.\ k \bmod 2$, and
\begin{align*}
  &\lemma{sum\_mod\_below}:\quad 0 < m\ \wedge\ \intf(n) \cdot (m - 1) \le e\\
  &\quad\Longrightarrow\ \mathsf{sum}(\mapc{\bl{\cnt(0, 1, n)}}{k \mapsto \ite{0 < k}{k \bmod m}{0}}) \le e\ \Longleftrightarrow\ \mathsf{true}
\end{align*}
closes the bound with $m := 2$, $e := N$ and $\intf(n) = N$ (61 rule applications). The three tests
\texttt{brs2}, \texttt{conda} and \texttt{ms2} share every step but the last two: the shape of the summands
(constants by index, constants by value, remainders) selects the introduction lemma for the stores and the
bound lemma for the sum.

\subsection{The idea}
Loops in sequence compose through the current values, as in the counting sort (Section~\ref{sec:cs}): the
stores of one loop are read in the update of the next, where the context of the iteration decides the
membership tests, and the read collapses to the stored value before the introduction lemma of the next loop
is tried. The order matters: \lemma{stores\_cases\_def} sees $a[k] + 1$ as $2$ because the read was resolved
first, which is why its first case is a constant; a lemma written from the program text, with $G'(a[k])$ in
the first case, would not match the simplified update. The lemmas for a program are written for the goal as
the simplifier prints it, not for the program (the pattern of Section~\ref{sec:principles:matcher}, arithmetic
in the goal's form).

\section{Strided Writes: Parity, Division, and a Case the Simplifier Does Not Close}\label{sec:pr2}

\subsection{The program}
The test \texttt{forall/pr2} fills an array of even size pairwise with values that depend on a bound
\texttt{MINVAL}, which the loop does not change, and asserts that every cell is at least the bound or zero.
\begin{lstlisting}[language=C]
if (CELLCOUNT % 2 != 0) return 1;
for (i = 1; i <= CELLCOUNT / 2; i++) {
  volArray[i*2 - 2] = (3 >= MINVAL) ? 3 : 0;
  volArray[i*2 - 1] = (1 >= MINVAL) ? 1 : 0;
}
for (i = 0; i < CELLCOUNT; i++)
  __VERIFIER_assert(volArray[i] >= MINVAL || volArray[i] == 0);
\end{lstlisting}
The run is \texttt{-unbounded -teq -simplify -split -simplify -include teq:strides -include
benchmarks/monocera/lemmas.dfy}: the module \texttt{strides} holds the lemmas on strided stores, and
\texttt{lemmas.dfy} the one lemma that is not in the database, \lemma{half\_even}.

\subsection{The trace equation system}
The loop runs from $i = 1$, so its trace of $i$ is $\cnt(1, 1, n)$ and the stores are at $2k$ and $2k + 1$
for the iteration $k = i - 1$. The body has four paths, two conditionals in sequence, of which the
simplifier's context removes the impossible ones (\texttt{MINVAL} $\le 1$ implies \texttt{MINVAL} $\le 3$),
and the merged update has three cases on the bound:
\begin{align}
  V &= \mathsf{trace}(v_0, \mathsf{for}\ \bl{\cnt(0, 1, n)}, \bl{\cnt(1, 1, n)}\ \mathsf{do}\ v, k, i \mapsto {} \nonumber\\
  &\qquad \upd{M \le 1}{v[2k := 3][2k + 1 := 1]}\ \ \upd{1 < M \le 3}{v[2k := 3][2k + 1 := 0]} \nonumber\\
  &\qquad \upd{3 < M}{v[2k := 0][2k + 1 := 0]}) \label{eq:pr1}\\
  &\al{\bl{\cnt(0, 1, n)}}{k \mapsto k < C / 2} \qquad 1 + \intf(n) = i' \qquad \neg(i' \le C / 2) \label{eq:pr2}\\
  &\al{\bl{K'}, \bl{I'}}{k, i \mapsto v'[k] \ge M \vee v'[k] = 0} \qquad i'' = \intf(n') < C \label{eq:pr3}\\
  &\Longrightarrow\ v'[i''] \ge M \vee v'[i''] = 0 \nonumber
\end{align}
with $C$ for \texttt{CELLCOUNT}, $M$ for \texttt{MINVAL}, $v' = \cur(V)$, and $C \bmod 2 = 0$, $1 < C$ from
the harness. The constraint \eqref{eq:pr2} is resolved by \lemma{invs\_counter\_upper} to
$n = 0 \vee (0 < C/2 \wedge \intf(n) - 1 < C/2)$, and with the exit $C/2 < 1 + \intf(n)$ this is the fact
$\intf(n) = C / 2$: the number of iterations is the half of the size, as a quotient.

\subsection{The lemmas}
\begin{align*}
  &\lemma{select\_current\_trace\_stride2\_cases}:\quad R \Leftrightarrow (\neg P \wedge \neg Q)\\
  &\quad\Longrightarrow\ \cur(\mathsf{trace}(a_0, \mathsf{for}\ \bl{\cnt(0, 1, n)}, \bl{\cnt(1, 1, n)}\ \mathsf{do}\ a, k, i \mapsto {}\\
  &\qquad\qquad \upd{P}{a[2k := X_1][2k + 1 := Y_1]}\ \ \upd{Q}{a[2k := X_2][2k + 1 := Y_2]}\\
  &\qquad\qquad \upd{R}{a[2k := X_3][2k + 1 := Y_3]}))[j]\\
  &\qquad = \mathsf{if}\ 0 \le j < 2\,\intf(n)\\
  &\qquad\quad \mathsf{then}\ (\mathsf{if}\ j \bmod 2 = 0\ \mathsf{then}\ (\ite{P}{X_1}{\ite{Q}{X_2}{X_3}})\\
  &\qquad\qquad\qquad\quad \mathsf{else}\ (\ite{P}{Y_1}{\ite{Q}{Y_2}{Y_3}}))\\
  &\qquad\quad \mathsf{else}\ a_0[j]\\[4pt]
  &\lemma{div\_nonpos}:\quad 0 < c\ \Longrightarrow\ (x / c \le 0) \Leftrightarrow (x < c)\\[4pt]
  &\lemma{half\_even}:\quad x \bmod 2 = 0\ \Longrightarrow\ (n = x / 2) \Leftrightarrow (x = 2n)
\end{align*}
The stride lemma reads a cell of the final array: the pairs $2k, 2k + 1$ are disjoint and cover
$0, \ldots, 2n - 1$, outside of which the initial array remains, and the value is chosen by the parity of
the cell and by the guards, which do not depend on the iteration. Its side condition requires the guards to
be exhaustive; without it the update would be unspecified where none holds and the equation unsound (the
guards of a trace need not be complete since \texttt{break}, Section~3 of the paper as implemented). The
lemma fixes two stores and three cases on its left-hand side, as the summaries of the algorithms fix the
fragment structure of one program; \texttt{pr3} to \texttt{pr5}, with $2^s$ combinations of $s$ per-cell
guards, have no lemma.

\subsection{The proof}
Of the 81 rule applications over the goals of \texttt{-split}:
\begin{enumerate}[nosep]
\item \textbf{The reads.} \lemma{select\_current\_trace\_stride2\_cases} fires four times, on the read $v'[k]$
  in the constraint \eqref{eq:pr3} of the checking loop and on $v'[i'']$ in the conclusion (each twice, in
  the two halves of the disjunction), with $P := M \le 1$, $Q := 1 < M \le 3$ and the six values; the
  side condition holds since the simplifier produced the three guards as a case split on $M$. Each read
  is now a conditional of constants below $2\,\intf(n)$, every branch of which satisfies
  $x \ge M \vee x = 0$ under its guard, so the assertion holds for $0 \le i'' < 2\,\intf(n)$.
\item \textbf{The harness.} \texttt{-split} separates the case of no iteration, $C / 2 \le 0$, in which the
  array is the initial one; \lemma{div\_nonpos} turns $C / 2 \le 0$ into $C < 2$, which contradicts
  $1 < C$, three times (one per goal that carries the case).
\item \textbf{The cells beyond the stores.} What remains is $2\,\intf(n) \le i'' < C$ with $\intf(n) = C / 2$
  and $C \bmod 2 = 0$: an empty range, which the simplifier does not see, since it inlines solutions of
  arithmetic equations into polynomials only and $C / 2$ is an atom to it, not a polynomial in $C$. Without
  \lemma{half\_even} this goal remains (the run with the module \texttt{strides} alone leaves exactly it,
  with the constraint of the checking loop fully resolved). \lemma{half\_even} rewrites the fact
  $\intf(n) = C / 2$ under $C \bmod 2 = 0$ into $C = 2\,\intf(n)$, a polynomial equation, which the
  simplifier solves and substitutes: $2\,\intf(n) \le i'' < 2\,\intf(n)$ is false.
\end{enumerate}
The constant strides without cases, \lemma{select\_current\_trace\_stride2} and \lemma{stride3}, close two
further tests in the same way and without the harness,\\
\texttt{array\_doub\_access\_init\_const} (stride 2) and \texttt{array\_tripl\_access\_init\_const} (stride 3).\\
In the first, $N = 100000$ is a literal, the
loop \texttt{i <= N} is a counter of $n$ steps with $\intf(n) = 100001$ decided from the constraint and the
exit, and the read $a'[k]$ for $k \le 200000$ is
\[ \ite{0 \le k < 2\,\intf(n)}{(\ite{k \bmod 2 = 0}{0}{0})}{a_0[k]} \]
(44 rule applications).

\subsection{The idea}
Three obstacles of different kinds in one small program. The stride is a shape of the update, answered by a
lemma that fixes the shape; the quotient in the bound is an obstacle of the simplifier's arithmetic,
answered by a lemma that restates the fact in the form the decision procedure handles (the pattern
``arithmetic in the goal's form'' of Section~\ref{sec:principles:matcher}, here the form the \emph{engine}
wants); and the exhaustiveness of the guards is a soundness premise that the lemma must state, since the
engine does not. \lemma{half\_even} stays outside the database as a workaround for a documented limitation
(TODO.md: solutions are inlined into polynomials only), the proper fix being in the simplifier.

\section{Stores Ahead: Upwards Proved, Downwards Diagnosed}\label{sec:rewnif}

\subsection{The programs}
The test \texttt{forall/rewnif} writes one cell ahead before it writes the current one; \texttt{rewnifrev}
(under \texttt{forall} and \texttt{min}) does the same in a loop that runs downwards. Both assert
\texttt{a[i] >= i} for all cells.
\begin{lstlisting}[language=C]
for (i = 0; i < SIZE; i++) { if (i + 1 < SIZE) a[i+1] = i; a[i] = i; }        // rewnif
for (i = SIZE - 1; i >= 0; i--) { if (i - 1 >= 0) a[i-1] = i - 2; a[i] = i; }  // rewnifrev
for (i = 0; i < SIZE; i++) __VERIFIER_assert(a[i] >= i);
\end{lstlisting}
The runs use the module \texttt{strides} (\texttt{-include teq:strides}), which holds the lemmas on stores
ahead. The summary of the MonoCera tests records \texttt{rewnif} as closed and \texttt{rewnifrev} as the one
test whose lemma ``has the same form as the trace in the goal when printed, yet does not match'', with the
literal $-1$ excluded as the cause. This section re-traces the first and resolves the second.

\subsection{Upwards: the trace equation system and the proof}
With $S$ for \texttt{SIZE}, $1 < S$ from the harness, the filling loop is
\begin{align}
  A &= \mathsf{trace}(a_0, \mathsf{for}\ \bl{K}, \bl{I}\ \mathsf{do}\ a, k, i \mapsto \upd{i + 1 < S}{a[i + 1 := i][i := i]}\ \ \upd{\neg(i + 1 < S)}{a[i := i]}) \label{eq:rw1}\\
  I &= \tr{0}{\bl{K}}{i, k \mapsto \upd{i + 1 < S}{i + 1}\ \ \upd{\neg(i + 1 < S)}{i + 1}} \nonumber\\
  &\al{\bl{K}, \bl{A}, \bl{I}}{k, a, i \mapsto i < S} \nonumber
\end{align}
and the lemma of the module is
\begin{align*}
  &\lemma{select\_current\_trace\_ahead}:\quad \intf(n) \le m\\
  &\quad\Longrightarrow\ \cur(\mathsf{trace}(a_0, \mathsf{for}\ \bl{\cnt(0, 1, n)}, \bl{\cnt(0, 1, n)}\ \mathsf{do}\ a, k, i \mapsto {}\\
  &\qquad\qquad \upd{k + 1 < m}{a[k + 1 := F(k)][k := G(k)]}\ \ \upd{m = k + 1}{a[k := G(k)]}))[j]\\
  &\qquad = \mathsf{if}\ 0 \le j < \intf(n)\ \mathsf{then}\ G(j)\\
  &\qquad\quad \mathsf{else\ if}\ j = \intf(n) \wedge 0 < j < m\ \mathsf{then}\ F(j - 1)\ \mathsf{else}\ a_0[j]
\end{align*}
The value written ahead at $k + 1$ is overwritten by the next iteration, so a cell $j < n$ holds what
iteration $j$ writes last, and the cell $n$ holds what the last iteration wrote ahead, if it did. The side
condition $\intf(n) \le m$ makes the two guards cover all iterations: for $k < n \le m$, either $k + 1 < m$
or $k + 1 = m$. The proof takes 27 rule applications: the counter $I$ is merged and recognized, the
constraint gives $\intf(n) \le S$, which is the side condition, and \lemma{select\_current\_trace\_ahead}
fires twice on the reads $a'[k]$ of the checking loop's constraint and $a'[i'']$ of the conclusion, with
$m := S$ and $F := G := \lambda z.\ z$ (the stored values $i$ abstracted over the bound variable), giving
$\ite{0 \le j < \intf(n)}{j}{\ite{j = \intf(n) \wedge 0 < j < S}{j - 1}{a_0[j]}}$; the exit
$i'' = \intf(n'')$ with $\intf(n'') = S = \intf(n)$ puts the conclusion's index in the first branch
(\lemma{integer\_injective}), and $j \ge j$ holds.

\subsection{Downwards: the goal that remained}
For \texttt{rewnifrev} the counter runs from $S - 1$ down, and the simplifier's goal, after 54 rule
applications and with the module's lemma for the downward loop present but silent, is
\begin{align}
  A &= \mathsf{trace}(a_0, \mathsf{for}\ \bl{K}, \bl{I}\ \mathsf{do}\ a, k, i \mapsto \upd{1 \le i}{a[i + (-1) := i + (-2)][i := i]}\ \ \upd{i < 1}{a[i := i]}) \label{eq:rw2}\\
  a' &= \cur(A) \nonumber\\
  I &= \cnt(\intf(n) + i', -1, n) \qquad 0 = i' + 1 \qquad 0 < \intf(n) + i' < 99999 \label{eq:rw3}\\
  &\mathsf{above\_iterations}(n'', \mapc{\bl{K''}}{k \mapsto a'[k] + 1}) \label{eq:rw4}\\
  &i'' = \intf(n'') \le \intf(n) + i' \nonumber\\
  &\Longrightarrow\ i'' \le a'[i''] \nonumber
\end{align}
Two things have happened to the facts. The exit $\neg(0 \le i')$ with the constraint $0 \le i$ gives
$i' = -1$, and the simplifier has eliminated $S$ in favour of the counter: $S - 1 = \intf(n) + i'$, so the
start of $I$ is the term $\intf(n) + i'$ (with $i' = -1$ kept as an equation, not substituted into it). And
the constraint of the checking loop, $\al{\ldots}{k \mapsto k \le a'[k]}$, has been abstracted by
\lemma{above\_iterations\_def} of the database into the atom \eqref{eq:rw4} before the read $a'[k]$ was
resolved, since no lemma resolved it: the fact is hidden from any later rule on $a'[k]$ (finding~3 of the
algorithms, the lemma that ``swallows'' $k < t$ constraints). In \texttt{rewnif}, the read was rewritten
first and the constraint resolved by itself.

\subsection{The diagnosis}
The module's lemma is \lemma{select\_current\_trace\_ahead} mirrored: the dependencies $\bl{\cnt(0, 1, n)}$,
$\bl{\cnt(s, -1, n)}$, the cases $\upd{1 \le i}{a[i + (-1) := F(i)][i := G(i)]}$ and $\upd{i < 1}{a[i := G(i)]}$,
and the right-hand side $G(j)$ for $s - \intf(n) < j \le s$, $F(j + 1)$ at $j = s - \intf(n)$ if $0 \le j$ and
$0 < n$, else $a_0[j]$. Four variants were run on the test with the lemma injected as a user file:
\begin{center}\footnotesize
\begin{tabular}{@{}llll@{}}
\toprule
variant & step of the counter & position of the store ahead & result \\
\midrule
as in the module & literal $-1$ & $i + (-1)$ & does not fire \\
a & literal $-1$ & $i - 1$ & does not fire \\
b & variable $d$, condition $d = -1$ & $i + (-1)$ & does not fire \\
e & literal $-1$ & $P(i)$, pinned by $\forall i.\ P(i) = i - 1$ & does not fire \\
d & variable $d$, condition $d = -1$ & $P(i)$, pinned & \textbf{fires twice, closed} \\
\bottomrule
\end{tabular}
\end{center}
Both changes are needed, and both are patterns of Section~\ref{sec:principles:matcher} that the algorithms
had established: a literal $-1$ in a counter pattern is not matched (principle~8 of \texttt{NOTES.md},
\texttt{reverse.dfy}, whose counter down has its step as a variable with $d = -1$ as a side condition), and
arithmetic on a bound variable in a left-hand side, here the position $i - 1$ of the store, is not matched
and is pinned (Section~\ref{sec:merge}, the swaps at $j - 1$ of the insertion sort). The summary had excluded
the literal $-1$ because the database matches $\cnt(i, -1, n)$ elsewhere; it does so as a term, not as a
dependency pattern matched through the context of the goal, which is where it stands here. With variant d the
test closes with $s := \intf(n) + i'$, $F := \lambda z.\ z - 2$, $G := \lambda z.\ z$, also under
\texttt{min}; a false assertion \texttt{a[i] >= i + 1} leaves a goal. The variant does not cover
\texttt{rewnifrev2}, which stores ahead \emph{upwards} ($a[i + 1]$) in a loop that runs downwards, a third
shape, whose cell $j$ holds the value stored ahead by the iteration below it.

\subsection{The idea}
The database lemma is right and the engine applies it to the upward loop; the downward loop fails on two
literal patterns whose workarounds the algorithms had found in another guise. The lesson for the database is
that lemma statements must follow the matcher's patterns as much as the goal's forms: steps of counters as
variables with side conditions, positions computed from the bound variable as pinned higher-order
variables. The lesson for the engine is the same as finding~12 and finding~34 of the algorithms: a match that
fails without a visible reason costs more than its cause deserves, and a flag that reports which part of a
left-hand side did not match would have answered this in one run. The restated lemma is a candidate for the
module \texttt{strides}, which this document does not change.

\section{A Recurrence on the Previous Cell (open)}\label{sec:ifeqn2}

\subsection{The program}
The test \texttt{forall/ifeqn2}: two arrays filled by recurrences on the previous cell, and the closed form
asserted for every position.
\begin{lstlisting}[language=C]
for (i = 0; i < N; i++) { if (i == 0) a[0] = 4; else a[i] = a[i-1] + 4; }
for (i = 0; i < N; i++) { if (i == 0) b[0] = 1; else b[i] = b[i-1] + a[i-1]; }
for (i = 0; i < N; i++) __VERIFIER_assert(b[i] == 2*i*i + 2*i + 1);
\end{lstlisting}

\subsection{The goal that remains}
After 127 rule applications, all on counters and their reads, the goal is
\begin{align}
  A &= \tr{a_0}{\bl{K}, \bl{I}}{a, k, i \mapsto \upd{k = 0}{a[0 := 4]}\ \ \upd{0 < k}{a[k := a[k - 1] + 4]}} \label{eq:ie1}\\
  a' &= \cur(A) \nonumber\\
  B &= \tr{b_0}{\bl{K'}, \bl{I'}}{b, k, i \mapsto \upd{k = 0}{b[0 := 1]}\ \ \upd{0 < k}{b[k := a'[k - 1] + b[k - 1]]}} \label{eq:ie2}\\
  b' &= \cur(B) \nonumber\\
  &\al{\bl{K''}, \bl{I''}}{k, i \mapsto b'[k] = 2k^2 + 2k + 1} \label{eq:ie3}\\
  &i'' = \intf(n'') < N \nonumber\\
  &\Longrightarrow\ b'[i''] = 2\,i''^2 + 2\,i'' + 1 \nonumber
\end{align}
with $K = I = K' = I' = \cnt(0, 1, n)$, $\intf(n) = N$, and $K'' = I'' = \cnt(0, 1, n'')$ the counter of the
checking loop. The two filling loops are not resolved: a store at $k$ of a value read at $k - 1$ is not a
$\mathsf{stores}$ of a map (the value depends on the trace, not on the initial array), and no lemma of the
database names such a recurrence. The checking loop has become the constraint \eqref{eq:ie3}, the assertion
for the earlier iterations, and the goal is the assertion at the current index: an induction step without a
closed form of $b'$.

\subsection{What is missing}
This is the first of the five classes of open goals that the summary of the MonoCera tests lists
(recurrences on earlier entries with a polynomial closed form: \texttt{ifeqn1} to \texttt{ifeqn5},
\texttt{ifcomp}, \texttt{pcomp}, \texttt{sqm}, \texttt{nsqm}, \texttt{condm}), and TODO.md records that
even a lemma written for the shape,
\[ \cur(\tr{a_0}{\bl{\cnt(0, 1, n)}}{a, k \mapsto a[k := a[k - 1]]}) = \ldots, \]
does not match its instance: $k - 1$ as a store index is arithmetic on a bound variable, the same obstacle
as the swaps at $j - 1$ of the insertion sort, where the pin of Section~\ref{sec:merge} is the way out
(a higher-order $P(k)$ for the position, pinned to $k - 1$). What a proof would need beyond the match:
\begin{itemize}[nosep]
\item an introduction lemma that names the recurrence, e.g.\ $\mathsf{recur}(a_0, c, n)$ for
  $a[k] = a[k - 1] + c$ from $a[0] = c$, with the acceleration $\mathsf{recur}(\ldots)[j] = c\,(j + 1)$ for
  $0 \le j < n$: the closed form of an arithmetic progression, which is what $a'[k] = 4k + 4$ needs;
\item the same for $b$, whose increment is the entry of another named trace, i.e.\ a sum of the
  progression, $b'[k] = 1 + \sum_{j < k} (4j + 4) = 2k^2 + 2k + 1$: either a lemma on sums of affine maps
  (the database has $\mathsf{sum}$ over counters, not over maps of a counter with a polynomial body), or a
  recurrence solver as a lemma generator for the arithmetic fragment, which the ideas on the paper
  propose as related work (compositional recurrence analysis, closed forms for numerical loops);
\item alternatively, the checking loop as an induction: \eqref{eq:ie3} is the invariant at the earlier
  positions, and the goal at $i''$ follows from $b'[i''] = a'[i'' - 1] + b'[i'' - 1]$ and the
  constraint at $i'' - 1$, an instantiation at a witness with the recurrence unfolded once, which no rule of
  the database performs (the step from a constraint to a position is the pattern of
  Section~\ref{sec:lsearch}, but here the position's value is given by the recurrence, not by a named
  function).
\end{itemize}
Neither Cuvée nor MonoCera proves this test (nor TriCera, CPAchecker or SeaHorn in the paper's evaluation):
the instrumentation of the quantified assertion by a ghost variable still needs the polynomial invariant of
the filling loops. A lemma-based proof needs the closed forms, which is the representation that trace
equation systems are meant to provide, and a database that does not yet have them.

\section{Principles}\label{sec:principles}

This section collects what the examples have in common: the shape that every proof takes, the rules
that the matcher imposes on a lemma written for a program, the discipline that keeps such lemmas sound and
the rewriting terminating, the diagnostics that located each obstacle, and what the same principles look
like on the C programs of the MonoCera benchmarks (\texttt{benchmarks/monocera}). The sources are the numbered
findings of \texttt{NOTES.md}, the experiment threads kept in \texttt{pending/} (preludes, programs with wrong
postconditions, scripts and logs, Oct 10 and 11, 2026), \texttt{doc/teq.md}, and the summary and case studies
of the MonoCera tests. The state of the folder on Oct 12, 2026: 27 of its 30 files verify with
\texttt{-teq -simplify} (nine of them with \texttt{-include teq:nested}); open are \texttt{is\_sorted.dfy}
(one substitution of the exit into an argument of $\sorted$), the existential postcondition of
\texttt{max\_element.dfy} (no witness), and the lower bound (Section~\ref{sec:lb}, blocked by the engine).

\subsection{The shape of a proof}\label{sec:principles:shape}

Every verified example is closed by lemmas of the following kinds, in this order from generic to specific.
The first two kinds are the database's, the others are the algorithm's knowledge, stated at the beginning of
its file.
\begin{enumerate}[nosep]
\item \textbf{The database resolves single loops.} Counters, sums, stores, selects, counts, minima and
  maxima are introduced and accelerated by the modules of \texttt{-teq} without anything written for the
  program (the four loops of the counting sort up to the table of start positions, Section~\ref{sec:cs}).
\item \textbf{Fragment lemmas resolve inner counters.} Theorem~4.8 of the paper turns the trace equation
  system of an inner loop into fresh sequences inside the iteration constraint of the outer loop, which the
  simplifier distributes into fragments over incomparable dependency lists that do not see each other's
  facts (the two minimal problems, \texttt{invs\_counter\_equation.dfy} and \texttt{invs\_fragments.dfy}). The
  module \texttt{nested} does the resolution by rewriting, one fragment at a time in the order in which
  \texttt{-teq} produces them: the inner trace is a counter, its final value is its length, the test bounds
  the length, the exit fixes it, the numbers of inner iterations are a constant sequence
  (Sections~\ref{sec:nsum} and~\ref{sec:sel:frag}). The same lemmas serve every nested example.
\item \textbf{A summary states what an inner loop computes.} Its left-hand side is the outer trace, read
  through the current value, or the whole atom of the postcondition; its side conditions are the fragments,
  matched against the facts; its right-hand side is a trace of an abstract step, $\argmin$ swaps,
  $\mathsf{insert}$, $\mathsf{shake}$, $\neg\mathsf{nodup}$ (Sections~\ref{sec:sel:sum}, \ref{sec:cock},
  \ref{sec:fd}). An inner result stored at the outer position needs no summary but a chain of specific
  substitutions (Section~\ref{sec:rows}).
\item \textbf{An induction lemma states the invariant of the outer loop} as a quantified side condition over
  the abstract step, which the simplifier decides by step lemmas on atoms; the invariant is a predicate
  indexed by the iteration ($\mathsf{bubbled}$, $\mathsf{shaken}$, the sorted prefix), and the lemma
  constrains all steps of the trace (Sections~\ref{sec:sel:ind} and~\ref{sec:bs}). This is the scheme of
  \lemma{current\_trace\_bounded} of the database: no invariant is inferred, its shape is the lemma's
  knowledge, and the database proves the step.
\item \textbf{Rules on the atoms of the postcondition} characterize a named result: the bounds of
  $\argmin$ and $\mathsf{lb}$, the regions of a partition, the counts of $\mathsf{merged}$, the two cases of
  the exit for $\mathsf{distincts}$. The postcondition is never rewritten as a whole; each atom finds its
  rule, with the exit of the loop among the facts.
\item \textbf{Name the normal form.} The counter of a loop is not inlined into the facts, so a count over
  $\bl{I}$ and over $\bl{\cnt(0, 1, n)}$ differ (\texttt{count\_occurrences.dfy}); a lemma that names the
  term makes both sides equal. Quantified preconditions are stated as atoms of the specification
  ($\mathsf{keys\_below}$, Section~\ref{sec:cs}), postconditions in the database's vocabulary
  ($\mathsf{occurs}$, $\mathsf{distincts}$, $\mathsf{count}$ over $\mathsf{selects}$) where the standard
  form has no witness or no matchable shape.
\end{enumerate}
A data-dependent number of iterations (merge, gnome sort, the early exit of the first duplicate) is no
obstacle to this shape as long as the property can be stated on the final values and the exit is read from
the facts; what the paper's loop shapes exclude is the characterization of the traces, not the postcondition
(Section~\ref{sec:gnome}).

\subsection{Writing a lemma that the matcher accepts}\label{sec:principles:matcher}

The matcher of the paper (Section~5.2) relates the binding structure of combinators, follows the definitions
of the context on demand, synthesizes dependencies that are variables of the rule as maps, and matches side
conditions against the facts while they have an open variable. From these four mechanisms follow the
patterns below, each found by an example that did not close without it (Table~\ref{tab:principles}).

\begin{table}[t]
\centering\small
\begin{tabular}{@{}lll@{}}
\toprule
pattern & found on & finding \\
\midrule
specific left-hand side of an $\mathsf{all}$ rule & cocktail sort, nested mult & 15, 29 \\
fresh variables for every $\mathsf{all}$ side condition & bubble sort, cocktail sort & 9, 27 \\
open trace variables for the traces of the goal & merge, counting sort, first duplicate & 14, 31 \\
the pin for literal patterns & merge, insertion sort, gnome sort & 14 \\
the whole atom as left-hand side & merge, gnome sort & 16 \\
arithmetic in the goal's form & cocktail sort, bubble sort & 11, 28 \\
side conditions in the forms of the included modules & insertion sorts, cocktail sort & 22, 37 \\
no condition on the rule's own term & lower bound, bubble sort & 20, 32 \\
the exit of an outer loop as two implications & first duplicate & 31 \\
atoms for constraints and quantified hypotheses & linear search, counting sort & 13, 19 \\
variables of the method named apart & fragments, invs\_fragments & 9, 21 \\
\bottomrule
\end{tabular}
\caption{The patterns of Section~\ref{sec:principles:matcher}, the examples that forced them, and the
findings of \texttt{NOTES.md} that record them.}\label{tab:principles}
\end{table}

\begin{enumerate}[nosep]
\item \textbf{A specific left-hand side.} An $\mathsf{all}$ pattern whose body has a bare bound variable of
  a variable dependency on one side, $n \le j$ or $j = \intf(nn)$, matches every comparison, respectively
  equation, fragment: the dependency is synthesized as a map of the fact's dependencies. With an
  $\mathsf{all}$ side condition such a rule cycles (Section~\ref{sec:principles:termination}). A function
  symbol on the side of the bound variable, $\intf(nn) \le n$ or $J = \cnt(0, 1, nn)$, an $\mathsf{all}$ with
  a specific body, or an atom of the specification is specific.
\item \textbf{The whole atom as left-hand side.} A variable argument of an uninterpreted function is not
  expanded by its definition when the simplifier rewrites below it: $\sorted(c, 0, n)$ with $c = \cur(C)$
  is not reached by a rule on $\cur(\tr{\ldots}{\ldots}{\ldots})$, but
  $\sorted(\cur(\tr{\ldots}{\ldots}{\ldots}), 0, s)$ is, and the length is a variable $s$ with $s = m + n$
  as a side condition, since a pattern $m + n$ of two variables is not matched (Section~\ref{sec:merge}).
\item \textbf{Fresh variables for every $\mathsf{all}$ side condition.} A condition is matched against the
  facts, with backtracking, while a variable of it is open; once all are bound it is closed and decided by
  the simplifier, which does not find a closed $\mathsf{all}$ among the facts (these carry unused
  dependencies, which the simplifier also drops from the instance, and their bound variables are renamed).
  So each $\mathsf{all}$ condition binds a fresh sequence, $\al{ks, js_6, Js_6}{k, j, J \mapsto j = \cur(J)}$,
  and an equation $js_6 = js_4 \wedge Js_6 = Js_3$ identifies it afterwards; the equation is closed and
  trivially decided. The 21st condition of the cocktail summary stayed closed for a day
  (Section~\ref{sec:cock}).
\item \textbf{Open trace variables.} The traces of the goal are variables $K$, $I$, $R$ with their equations
  as hypotheses. $\bl{\cnt(0, 1, m)}$ as a \emph{dependency} of a pattern matches $\bl{K}$ through the
  context; as an \emph{argument} of a function it never does, and a closed equation $is = \bl{\tr{\ldots}{\ldots}{\ldots}}$
  is decided and fails, since $\bl{I}$ is not expanded. A rule takes $Kv = \cnt(0, 1, m)$ or
  $Iv = \tr{0}{\ldots}{\ldots}$ as an \emph{open} condition, matched against the hypothesis, and uses
  $\bl{Kv}$, $\bl{Iv}$ afterwards (Sections~\ref{sec:merge}, \ref{sec:cs}, \ref{sec:fd}).
\item \textbf{The pin.} A literal pattern that the matcher does not handle, a disjunctive guard, the
  position $j - 1$ of a store, the initial value $n - k - 2$ of a counter, the step $-1$, becomes a
  higher-order variable $G(i, j)$, $P(j)$, $X(k)$, $d$, which matches anything, and a quantified side
  condition $\forall i, j.\ G(i, j) \Leftrightarrow \varphi$ or an equation pins it to the intended term; the
  simplifier decides the pin by beta reduction. A pin must mention only variables that the left-hand side
  or an earlier condition binds, and it must follow the fragments: a test $0 < j \wedge y[j] < y[j - 1]$ is
  split into two constraints, so its pin is two pins (\texttt{insertion\_sort.dfy}). A higher-order variable
  for a bare array read $a[i]$ is instantiated by the array itself (finding~10), so reads are pinned, as
  $U(i, j) = a[i]$ in the merge, or written literally.
\item \textbf{Arithmetic in the goal's form.} The decision procedure of the context combines difference
  constraints; $len \le 2\,\intf(m) + 1$ is not derived from $n \le \intf(m) + lo + 1$ and $lo = \intf(m)$.
  A condition states the bound as the goal has it, with the program's variable as an open rule variable
  matched against its equation ($l = \intf(m)$, then $len \le \intf(m) + l + 1$); a length $len - k - 1$
  that a pattern cannot contain is a variable $p$ with $p = len - k - 1$ as a closed condition, decided once
  $p$ is bound (Section~\ref{sec:bs}). Arithmetic on a bound variable is admissible on right-hand sides and
  in side conditions, never in a left-hand side.
\item \textbf{Side conditions in the forms of the included modules.} The fragments are rewritten by the
  modules before a summary sees them, so a summary is written for the normal form of the whole database
  that the run includes: $j = \cur(J)$ without \texttt{teq:nested}, $j = \intf(nn) + k + 1$ with it, as the
  goal prints them. A file includes exactly the modules its lemmas are written for; all nested loops of the
  folder are now written for \texttt{teq:nested}, so that the module can become part of \texttt{-teq}
  (finding~37).
\item \textbf{No condition on the rule's own term.} The instance of a side condition is simplified before it
  is decided or matched, so a condition $l = \cur(Lv)$ of a rule on $\cur(Lv)$ with $Lv$ bound rewrites
  $\cur(Lv)$ by the rule itself and does not terminate. The summaries take $\cur(\mathsf{results}(\ldots))$
  as left-hand side rather than $\mathsf{results}(\ldots)$ for the same reason, and the exit of a loop is
  read from the facts ($x = \cur(Iv)$ open, then $n \le x$, Section~\ref{sec:gnome}) or abstracted by an
  atom, never stated on the current value of the trace being rewritten (Section~\ref{sec:lb}).
\item \textbf{The exit of an outer loop.} The exit of an inner loop is a fragment and a side condition like
  the others. The exit of the outer loop stands in the conclusion of the goal, $\neg exit \vee post$, which
  \texttt{-split} does not split when the exit is a disjunction, and a rule on an atom of $post$ sees
  neither literal; a postcondition stated as two implications gives each goal its literal
  (Section~\ref{sec:fd}). The case of no iteration appears only after a first simplification, hence
  \texttt{-simplify -split -simplify} where it matters (\texttt{is\_sorted.dfy}, \texttt{reverse.dfy}).
\item \textbf{Atoms for iteration constraints and quantified hypotheses.} An iteration constraint is never a
  side condition (Section~\ref{sec:principles:termination}); a definition lemma abstracts it by an atom
  with a body that the matcher cannot synthesize, an explicit read $\mathsf{earlier}(\natf(k), js) = c$, and
  the lemma that needs the constraint takes the atom (Section~\ref{sec:lsearch}). A quantified hypothesis is
  matched by no side condition, open or closed; the specification provides the atom (Section~\ref{sec:cs}).
\item \textbf{Variables named apart.} Rule variables are not renamed apart from the variables of the goal.
  A method variable with the name of a rule variable of an open condition makes the condition coincide
  with the fact literally; it is decided true without binding the variable, and the run ends with
  ``variables of a rule not instantiated by its match'' on stderr, which \emph{replaces the goal}: zero
  goals with an error is not a proof. The fragments of \texttt{-teq} are named $K?_{14}$ and never clash;
  test methods and pins name their variables $KK$, $NN$, $len$ apart.
\item \textbf{Dependencies.} Lemmas carry the extra dependencies that the simplifier keeps in maps (the
  counter of the loop as $hs$), the matcher does not backtrack over the selection of dependencies when a
  side condition fails (hence \lemma{invs\_current\_counter\_other}), equal counters $K = I$ of one loop are
  not identified, and a higher-order summand $X(k)$ matches a single summand only, which is why the
  resolved forms are written literally.
\end{enumerate}

\subsection{Soundness}\label{sec:principles:sound}

The lemmas of the examples are axioms: the abstract predicates $\sorted$, $\perm$, $\partitioned$,
$\mathsf{bubbled}$, $\mathsf{shaken}$ and the functions $\argmin$, $\mathsf{insert}$, $\mathsf{shake}$,
$\mathsf{nodup}$, $\mathsf{merged}$, $\mathsf{lb}$ have no definitions, and the summaries state the
correctness of an inner loop in one line. Nothing in the tool checks them, so the discipline of the folder
replaces the check.
\begin{itemize}[nosep]
\item \textbf{Every result is checked on a wrong postcondition.} Each thread of \texttt{pending/} has a
  \texttt{program\_wrong.dfy} ($\sorted(b, 0, n + 1)$, the sum plus one, the negated flag); a prelude that
  closes its goal is unsound or vacuous, not a success. Two such cases occurred: a rule
  $\mathsf{nodup}(a, k, n) \Leftrightarrow \mathsf{true}$ under the hypothesis $\mathsf{distincts}(\ldots)$
  fired inside the context in which a side condition of a rule \emph{on} $\mathsf{distincts}$ was matched, so
  the atom was derived from itself (finding~13); and a condition with an open dependency and an equation on
  its bound variable was decided vacuously (\texttt{vac.dfy}). The lesson: a side condition must not be
  simplifiable by rules that use the conclusion, and the safe form of a condition is an atom that is matched
  literally.
\item \textbf{A lemma must be a lemma.} The lower bound closes with an exit $h \le l$ of two open variables
  matched against the fact $hi \le lo$ (\texttt{lb\_cmp0.dfy}); in this goal its conclusion is right, since
  the exit holds in the final state, but $h$ and $l$ are not tied to the traces and the rule is not a
  theorem. Such a rule is not admitted, and the file stays open (Section~\ref{sec:lb}).
\item \textbf{An induction constrains all steps.} A quantified condition over $0 \le k < n$ with
  $n \le \intf(m)$ steps leaves the further steps unconstrained. The number of steps is fixed exactly
  ($n = \intf(m)$, a fact of the goal), or the invariant is clamped so that every step keeps it
  ($\mathsf{bubbled}$ for $k \ge len$), or the result holds for any number of iterations beyond a bound
  ($m + n \le \intf(nn)$ in the merge).
\item \textbf{Soundness premises are the matcher's facts.} The Lean proofs of the core modules showed the
  premises that the lemmas need and the statements omit: dependencies of equal length, indices within the
  lists, complete guards; the fallback encoding does not assume them. The modules \texttt{nested},
  \texttt{sorting} and \texttt{strides} have no Lean proofs yet, and the abstract predicates would need
  definitions there for the step lemmas to be provable; both are open items.
\item \textbf{A killed run prints no goal.} A runner that counts goals reports a timeout as a proof: the
  first ``closed'' runs of \texttt{row\_sums}, \texttt{transpose} and most variants of the lower bound were
  kills at 120 or 300 s (finding~18). Every status of the folder was re-established with the exit status
  read, and long runs use a private copy of the jar, since a rebuild during a run crashes the JVM
  (finding~33).
\end{itemize}

\subsection{Termination}\label{sec:principles:termination}

Rewriting cycles by one mechanism, recorded for \lemma{sum\_invs\_const} of the database and met four
times in the folder (findings~19, 24, 29): an open $\mathsf{all}$ side condition is matched against the
$\mathsf{all}$ facts, which builds the simplifier's context of each fact's combinator and rewrites the reads
$\mathsf{earlier}(\natf(k), \bl{X})$ of its dependencies there; if the left-hand side of the rule being
decided matches one of those rewritten reads, the rule recurs through the same fact one level deeper until
the step budget is exhausted. Two ingredients, both needed: an $\mathsf{all}$ condition over lists that the
left-hand side binds, and a left-hand side general enough to match inside that context. Hence the rule of
thumb for a rule with $\mathsf{all}$ conditions: its left-hand side must be a term that no rewriting of a
read can produce, an $\mathsf{all}$ fragment with a specific body or an atom of the specification, never a
bare read, comparison, equation, or $\cur$ of a named trace for which an $\mathsf{earlier}$ acceleration
lemma exists or may be added. The general substitution lemmas of the first version of \texttt{row\_sums.dfy}
(a higher-order body $t = F(k, J)$, a read $\mathsf{earlier}(\natf(i), ts)$) are exactly what this forbids
(Section~\ref{sec:rows}); the exit lemma of the module \texttt{nested} with the body $n \le j$ cycled with
each counter lemma and with none alone, found by pair runs, and is stated on the constraint fragment
$\intf(nn) \le n$ now. The order of the modules also decides which abstraction a fragment gets:
\lemma{consts\_prefix\_def} of \texttt{aggregation} abstracts the exit $j = n$ before \texttt{nested} sees it,
hence \lemma{consts\_prefix\_integer} with the atom as left-hand side (finding~23).

\subsection{Diagnostics}\label{sec:principles:diag}

The obstacles of this folder were located by four techniques, each worth a flag of the tool:
\begin{itemize}[nosep]
\item \textbf{Preludes.} A program without lemmas and a file of lemmas spliced in front of it by the
  script \texttt{with.sh} of \texttt{pending/}, so that variants of the lemmas run against the same program;
  the folder holds the series for the cocktail sort (\texttt{cock2} to \texttt{cock9b}), the first duplicate
  (\texttt{fd4} to \texttt{fd8}), the lower bound, the row sums and the transpose, and the bisection of the
  module \texttt{nested}.
\item \textbf{Truncated copies.} A summary with its first $k$ side conditions, $k$ downwards, finds the first
  failing condition (the 21st of the cocktail summary); the copies are unsound and serve the diagnosis only.
\item \textbf{Probes.} Rules with the real left-hand side, the first $k$ conditions and a fresh atom
  $\mathsf{probe}_k(\ldots)$ as right-hand side show in the goal the largest $k$ whose conditions hold, and a
  probe that fires with ``variables of a rule not instantiated'' shows a left-hand side that left variables
  unbound (finding~32). A flag \texttt{-probe <lemma>} that reports the first failing condition of a named
  rule, with the facts it was matched against, would replace an afternoon of probe files.
\item \textbf{Pair runs.} The single lemmas of a module in pairs on the program that cycles find the pair
  that does (\texttt{pending/nested\_module/pairs.sh}); the rule reported at the step limit is the last step
  of the loop, not its cause.
\end{itemize}
With \texttt{-rules} a proof reports the lemmas it used; a residual goal points to the missing one. The
whole folder runs in about five minutes on one core, 5 to 35 s per file, the row sums and the transpose in
under two minutes.

\subsection{The same principles on C programs: the MonoCera tests}\label{sec:principles:monocera}

The 125 SV-COMP tests of MonoCera (Amilon et al., CAV 2023; instrumentation of $\forall$, $\max$, $\min$ and
$\sum$ over arrays) are C programs with a loop that fills an array and a loop that checks an assertion over
it. Their summary (\texttt{benchmarks/monocera/summary.md}) and the case studies of the lemmas written for
them (\texttt{lemmas.tex} there, Cuvée 0.35.15) show the principles of this folder on another class of
programs, without hand-written specifications.

\begin{table}[t]
\centering\small
\begin{tabular}{@{}lrrrrr@{}}
\toprule
operator & tests & MonoCera & Cuvée \texttt{-teq} & Cuvée \texttt{-chc} & TriCera, CPAchecker, SeaHorn \\
\midrule
forall & 93 & 26 & 53 & 0 & 0, 1, 0 \\
max & 4 & 2 & 2 & 0 & 0, 0, 0 \\
min & 11 & 3 & 7 & 0 & 0, 0, 0 \\
sum & 17 & 10 & 17 & 0 & 0, 0, 0 \\
\midrule
total & 125 & 41 & 79 & 0 & 0, 1, 0 \\
\bottomrule
\end{tabular}
\caption{Proved tests of the MonoCera evaluation (Cuvée 0.40.0, 15 min, Oct 12, 2026; the other tools' numbers
from the paper). All proofs of \texttt{-teq} end within 20 s (median 9 s); no tool gives a wrong result.}
\label{tab:monocera}
\end{table}

\begin{itemize}[nosep]
\item \textbf{The shape of the data decides the lemma, not the idea.} An accumulator kept in the cell
  \texttt{sum[0]} of a one-element array is a trace of stores, which \lemma{current\_trace\_reset\_add} of the
  database, written for a scalar, does not match; \lemma{select\_current\_trace\_reset\_add} is the same lemma
  for the read of the cell, with the scalar lemma's right-hand side so that the database continues from
  there. This is ``name the normal form'' (Section~\ref{sec:principles:shape}): the Dafny translations of
  SV-COMP make the cell a scalar and close the test for free, the C programs need the variant.
\item \textbf{Left-hand sides bind, side conditions must be facts.} \lemma{sum\_cases\_below} bounds a sum of
  $n$ terms that are each $0$, $c$ or $d$ by $e$ with its left-hand side a comparison whose constants bind
  the bound; the general \lemma{sum\_below\_prefix} of the database does not apply, since its bound occurs
  only in a side condition, which would have to be among the facts. The counting sort's
  \lemma{rank\_table\_def} binds the key bound $kk$ for the same reason (Section~\ref{sec:cs}), and the
  pattern $s = m + n$ of the merge is its arithmetic form.
\item \textbf{Lemmas compute values, they do not make goals vacuous.} Each closed test was rerun with a
  false assertion, which stays open with a residual such as $\neg(0 < N)$: the lemma computed $N - 1$ and the
  bound was wrong. The same check as the wrong postconditions of Section~\ref{sec:principles:sound}.
\item \textbf{Specific lemmas are summaries.} The strided stores (\lemma{select\_current\_trace\_stride2},
  stride 3, three exhaustive cases at stride 2) and the stores ahead ($F(k)$ at $k + 1$, then $G(k)$ at $k$)
  fix the number of stores and of cases on their left-hand sides, as the summaries of this folder fix the
  fragment structure of one algorithm; strides 3 to 5 with per-cell cases, writes by inner loops and
  recurrences on earlier entries remain open. Both want what the arity variants \lemma{perm\_steps},
  \lemma{perm\_steps\_two} and \lemma{sorted\_prefix\_steps\_two} want: a pattern with the stride, the number
  of dependencies, or the number of cases as a parameter.
\item \textbf{The checking loop is an instantiation at a witness.} In 72 of the 96 open goals of the
  earlier run the assertion of the checking loop survives as an $\mathsf{all}$ over its earlier iterations,
  and the goal is the assertion at the current index $i = \intf(n)$ of a filled array whose closed form is
  not derived (recurrences $b[i] = b[i - 1] + a[i - 1]$, parity, division): the step from the constraint of
  one loop to a position of another is the one that \lemma{not\_occurs\_selects\_def} and
  \lemma{select\_not\_occurs\_selects} make for the linear search (Section~\ref{sec:lsearch}), an atom and
  its instantiation, and the missing closed forms are the partial descriptions of Section~7 of the paper
  (the $\max$ and $\min$ tests, the lower bound here).
\item \textbf{The higher-order variable and the array.} The simplifier's eta reduction turned the instance
  $(z) \mapsto a[z]$ of a higher-order summand $G$ into the array $a$, leaving an application $a(k)$ that no
  rule about $a[k]$ matches; restricting eta reduction to lambdas of array type fixed it in 0.35.15. The
  examples of this folder meet the same phenomenon in the matcher (finding~10) and avoid a higher-order
  variable for a bare read, reading the array literally (\texttt{prefix\_sums.dfy}) or pinning the stored
  values (Section~\ref{sec:merge}).
\item \textbf{A match that fails for no visible reason is usually one of the known patterns.} The lemma for
  the stores ahead of a loop that runs downwards had the same printed form as the trace in the goal and did
  not match, and the summary of the MonoCera tests had excluded the literal $-1$ as the cause. Re-traced here
  (Section~\ref{sec:rewnif}), it is two of the patterns of Section~\ref{sec:principles:matcher} at once: the
  literal step $-1$ of the counter dependency is not matched (a variable $d$ with $d = -1$ as a side
  condition is, as in \texttt{reverse.dfy}), and the position $i - 1$ of the store ahead is arithmetic on a
  bound variable (the pin). With both, the lemma fires and the two \texttt{rewnifrev} tests close
  (\texttt{rewnifrev2} stores ahead upwards in a loop that runs downwards, another shape). The cocktail
  summary that fired once in a day (finding~12) and the lemma of the module \texttt{nested} that fires in
  the input but not from the module (finding~34) are residues of the same kind, for which the probe flag of
  Section~\ref{sec:principles:diag} is the remedy.
\item \textbf{Lookup against inference.} \texttt{-chc} with the Horn solver of z3 proves none of the 125
  tests (120 timeouts): the quantified invariants over arrays that the tests need, $b[k] = k^2 + k + 1$ for
  all $k < i$, are not found, as for TriCera in the paper, and the 15 minutes add nothing to either
  MonoCera or \texttt{-teq}, whose proofs end within 20 s. This is the paper's thesis in numbers: a lemma
  found once serves every loop of its shape, and the price is that a loop without a lemma stays open
  (43 of the 125, mostly the recurrences, the parity and the nondeterministic input arrays).
\end{itemize}

\subsection{What the engine and the database need}\label{sec:principles:next}

The workarounds of Section~\ref{sec:principles:matcher} are each a missing feature of the engine, ordered
by the payoff that the examples measure: a re-entrance guard, such that a rule is not applied while the
instances of its own side conditions are simplified (the sound lower bound closes, the summaries lose the
$\cur(\mathsf{results}(\ldots))$ detour); closed $\mathsf{all}$ conditions matched modulo unused
dependencies instead of decided (the fresh variables and identifying equations disappear from five files,
the summaries shrink by a third); a left-hand side matched through the definition of a trace variable binds
all its variables or is refused; a static check that rejects a left-hand side with a bare bound variable of a
variable dependency together with an $\mathsf{all}$ condition (it would have refused the cycling exit lemma
on load); and a case split of \texttt{-split} on a disjunctive negated exit. For the database: the module
\texttt{nested} into \texttt{-teq} with a run on the SV-COMP translations; the chain of the row sums
generalized into it (counters over constant numbers of iterations, final values given per iteration,
pointwise definitions as maps, reads of a map over the counter); the module \texttt{sorting}, which no file
includes yet, taken up by the examples; Lean proofs for \texttt{nested}, \texttt{sorting} and
\texttt{strides}, with definitions of the abstract predicates; and lemma schemas, one per observer of a
named scan, in place of the arity and type variants. The algorithms that remain open need, respectively,
the substitution of a solved exit into a function argument (\texttt{is\_sorted.dfy}), a witness for an
existential postcondition or quantified left-hand sides (\texttt{max\_element.dfy}), and the re-entrance
guard (the lower bound).

\section{Conclusion: Symptoms, an Engine, and the Representation}\label{sec:conclusion}

\subsection{The patterns are symptoms}
Section~\ref{sec:principles} presents the twelve patterns of Table~\ref{tab:principles} as a discipline for
writing lemmas. Read against their causes, they are a bill for the engine. Nine of the twelve are the matcher
(specific left-hand sides, fresh variables for closed conditions, the pin, variables named apart), the context
(open trace variables, the whole atom as left-hand side, arithmetic in the goal's form), or the rewriting
strategy (side conditions in the forms of the included modules, no condition on the rule's own term, the exit
of an outer loop as two implications); one is the representation (the fragments of inner loops that do not
share facts); only the atoms for quantified hypotheses and iteration constraints are partly a design choice,
what a lemma may take as a fact. The residues of the MonoCera tests sort the same way: the downward stores-ahead
lemma failed on a literal step and a position computed from a bound variable (Section~\ref{sec:rewnif}),
\lemma{half\_even} exists for a solution that is not substituted (Section~\ref{sec:pr2}), and a definition lemma
hid a constraint by rule order. What the examples strained is not the method, lookup of loop shapes in a
database of lemmas, but the engine that performs the lookup: syntactic matching modulo a fixed list of
theories, first match without backtracking, a context that follows definitions only where the matcher asks,
conditions simplified recursively, and normal forms that depend on the order of the rules.

\subsection{Trace equation systems: an interface, not the sole carrier}
A loop over variables $\bar x$ is, semantically, one fold over the iterations,
$S = \tr{\bar s_0}{\bl{K}}{\bar s, k \mapsto F(\bar s, k)}$, with the state tuple as the accumulator. A trace
equation system is this fold decomposed by projection: one equation per variable, with the other components as
sequences among the dependencies. The decomposition is exact, every projection of a fold over a product is a
fold over the product's history, and symbolic execution computes it for free. Where the dependencies form a
layering (a counter, an index over it, an array over both, a sum over the array), the decomposition is the
compositional reasoning of the method: each equation meets a small lemma of the database and the results are
substituted downstream, which is why the single loops of this document and the arrays of SV-COMP go through
with nothing written for the program. A single fold over the state would first have to discover that layering
before any lemma on a counter or a sum could apply.

Where the graph of dependencies is not a layering, the decomposition is the obstacle, and these are the hard
examples. Mutually dependent variables, the pointers of the merge, the bounds of the binary search, the position
and the array of gnome sort, give equations that are individually meaningless, and every lemma reconstructs the
product by tying the dependency lists back to their traces with open trace variables. Inner loops become fresh
sequences distributed into fragments over incomparable dependency lists, which is the whole fragments problem of
Sections~\ref{sec:nsum} and~\ref{sec:sel:frag}. Equal counters that are not identified, the injective selection
of dependencies, unused dependencies dropped and later missed, and the arity variants of lemmas are bookkeeping
of dependency lists that a single shared dependency would not have.

The conclusion we draw is that trace equation systems are the right interface for the lemma database and the
wrong sole carrier of the goal. The state trace of a loop should be the primary object, with one counter, one
dependency list $\bl{K}, \bl{S}$, all iteration constraints over it, and the inner loop as a value of the update;
the per-variable equations are views, $X = \pi_x(S)$ with $\bl{Y} = \mapc{\bl{S}}{\bar s \mapsto \pi_y(\bar s)}$,
derived by one projection rule whenever the update of a component mentions only some components, which is
exactly when a lemma of the database can apply to it. The matcher already synthesizes dependencies as maps when
a pattern dependency is a variable; this is the same mechanism made principal. The lemmas keep their form.

\subsection{An engine without the symptoms}
Five components, from the most local to the most far-reaching, each removing a group of patterns
(\texttt{ideas/deduction-engine.md} has the details and an order of work):
\begin{enumerate}[nosep]
\item \textbf{Congruence closure as the context}, in place of the chains of variable equations and the
  definitions followed on demand: a pattern variable matches an equivalence class, a definition is followed in
  every position, both orientations of an equation are available, equal counters are identified, and solutions
  reach function arguments. This closes \texttt{is\_sorted}, retires \lemma{half\_even} and \lemma{mul\_eq\_cancel},
  and removes the open trace variables and the whole-atom left-hand sides.
\item \textbf{Canonical forms before matching}, with closed conditions looked up modulo those forms instead of
  re-simplified, backtracking over the selection of dependencies, rule variables in their own namespace, and
  abstraction applied by the engine where a pattern subterm over bound variables fails to match. No fresh
  variables with identifying equations, no pins, no literal $-1$, no arity variants.
\item \textbf{Saturation with triggers} in place of the recursive simplification of condition instances:
  conditional rules fire when their conditions are present, facts about terms that occur are derived forward
  (the bounds of named monotone traces, the body of an iteration constraint at an index in range, the value of a
  trace at its exit), and a disjunctive exit is split by the engine when a rule on an atom needs a literal of it.
  This removes the cycles, the condition on the rule's own term, the atoms for constraints, and the dependence
  of normal forms on the order of modules; a condition that is a formula to prove, the step of an induction
  lemma or a pin, is a subgoal with the parent rule suspended.
\item \textbf{Equality saturation for the terms}, the long-term form of the first three: rules add equal nodes
  rather than replacing, extraction picks the result, nested e-graphs per combinator body carry the iteration
  facts.
\item \textbf{State traces with projections}, as above, which dissolves the fragments problem and makes the
  lemmas for mutually dependent shapes what their summaries are today after all the tying: a lemma on a fold
  over a pair with the guard as a higher-order variable.
\end{enumerate}
With the first three and the representation, \texttt{is\_sorted} closes by congruence, the nested sums and the
two minimal problems need no fragment lemmas, the merge, gnome sort and the lower bound take one lemma on a fold
over a pair, the sorting summaries take the inner fold as one condition instead of seven to twenty-one, the first
duplicate keeps its equivalence, the existential postcondition of \texttt{max\_element} finds its witness among the
terms, and the stores-ahead lemma matches downwards as it does upwards. What stays is what the method is for: the
database of introduction and acceleration rules, the symbolic execution that produces the equations, and the
summaries and inductions that state what an algorithm's loops compute.


\section*{AI Declaration}

The theory of trace equation systems, their generation from loops, the rewriting strategy and the matching
of traces are the work of the author and his coauthors in \emph{Trace Equation Systems for Program
Verification}; the implementation in Cuvée is described in the competition contribution of Cuvée to
SV-COMP~2027 (\texttt{doc/paper}), with its own declaration. The examples of \texttt{examples/dafny/algorithms},
their specifications and lemmas, the modules \texttt{nested} and \texttt{sorting} of the lemma database, the
findings of \texttt{NOTES.md}, the experiments kept in \texttt{pending/}, the lemmas for the MonoCera tests and
this document were produced by coding agents (Claude Code with Claude Opus; the restructuring of the document,
Sections~\ref{sec:monocera} and~\ref{sec:principles} with Claude Fable~5.1) between October 9 and 12, 2026,
under the direction of the author, who chose the problems, set the rules of the work (no changes to the
engine, no unfolding of trace combinators in lemmas for programs, every result checked on a wrong
postcondition, runs with their exit status), decided the open questions, and reviewed the results. The runs
reported here were made by the agents with the tool; the lemmas of the examples are axioms that no proof
checks (Section~\ref{sec:principles:sound}), and the text should be read as the agents' account of the
experiments, reviewed by the author.

\end{document}
